% Modernised reading — fonds Alexandre Grothendieck, shelfmark 161-6
% (pages 1-33, the whole folder).
%
% This is an INTERPRETATION, not a transcription. It re-reads the pages in
% today's notation and vocabulary, states the hypotheses the manuscript leaves
% implicit, and drops the critical apparatus so that the mathematics reads
% straight through. Everything the brackets and underlines carried moves into a
% footnote. It is held to being mathematically correct as it stands; where the
% manuscript is loose or mistaken, the true statement is what appears here and
% the footnote says what the page has. In any dispute of substance the
% transcription governs: whoever cites Grothendieck must cite
% batch-01.fr.tex (pp. 1-20) or batch-02.fr.tex (pp. 21-33).
%
% Pass: Opus 5.5 (claude-opus-5-5), 2026-09-22 — first pass, unchecked against the pages by a human.
% The folder was read whole, both batches together; this is one document for
% the shelfmark. It was made on Opus 5.5 and has not been measured against the
% reference reading of folder 115 (made on Opus 5).
% Revised 2026-09-23 (Opus 5.5): no order in time asserted between these
% undated leaves and anything else — the résumé and the footnote to VII no
% longer say he « reviendra » to these objects in the Esquisse (1984) or call
% it « postérieur », ribbon/fat graph are dated « des années 1980 » without
% « postérieurs », 162-1 is no longer « de la même période », p. 17 is no
% longer a « premier jet » of pp. 3-4 nor corrected by the fair copy, the
% p. 3 margin is no longer the « germe » of p. 21, p. 21's golden constant is
% « une autre idée », not « une première idée »; p. 18 footnote: the alpha-chain
% is exact, the pointer to the (now corrected) transcription note removed;
% chemise name given as Labarde or Laborde, undecided (facsimile re-read);
% pencil run: only pp. 15 and 13 are struck, one oblique stroke each, the two
% lines at the foot of p. 12 are not, and only p. 13 and the foot of p. 12 are
% head to tail to the ink (spine list and section V), checked on the facsimile.
%
% Spine. Not one argument but three subjects and a digression, on loose leaves.
%  (1) The icosahedron: ink fair copy pp. 3-10 (construction, calottes, lemma,
%      Gamma' simply transitive on oriented edges, |Gamma'| = 60, the two
%      graphs, the wish to see Gamma' combinatorially); notebook leaves
%      pp. 16-19 (pp. 3-4's computation in other letters, Gamma' = A_5 via five triples, the coset
%      lemmas Hom_G(G/H,E) = E^H and Aut_G(G/H) = N(H)/H); pp. 14 then 12 (the
%      golden number in Z[1/5][zeta_5]); p. 21 (icosahedra relative to a
%      quadratic form over a base).
%  (2) Graphs: pp. 24-30 (isotopy classes of embeddings of a 1-complex in an
%      oriented surface = ribbon graphs, model surfaces, mapping class group,
%      normal and planar embeddings, a Question, a fibration sequence);
%      pp. 32-33 (graphs with a group transitive on vertices and oriented
%      edges, coded by G > H' > H and sigma-dot).
%  (3) Products: pencil run pp. 15 -> 13 -> foot of 12 (a coproduct in a
%      category « (OM) » of ordered sets, read here as frames); p. 22,
%      cancelled (pi(C,E) = Hom(C°,E) and gamma(C,E) : pi(C,E) -> C^ x E),
%      footnoted against 161-4 pp. 15, 17 (same notation, not merged).
%  The links between (1) and (2) — p. 10's coherent circular orders are a
%  rotation system; the icosahedron is the coset graph of (A_5, 1 < C_5,
%  sigma) of pp. 32-33 — and between the two halves of (3) are OURS, said so.
%
% Notation fixed for the document (each part keeps the page's letters except):
%  - rotation of angle t about a unit vector a: R_{a,t}; the page's
%    gamma = gamma^w_alpha gamma^u_{-theta/2} (pp. 5-6) is written r.
%  - eta = zeta + zeta^{-1} = 2cos(2pi/5), eta' = zeta^2 + zeta^{-2}; the page
%    writes gamma (pp. 12, 14, 21), eta (p. 14), t, c. The ring
%    Z[1/5][t]/(t^2+t-1) and its spectrum: theta_5 (vartheta), which p. 14
%    also writes V_n and p. 21 double-struck V_5.
%  - (OM) run: m : I x J -> K for the page's gamma, a, b for its alpha, beta,
%    f_1, f_2, f for its alpha-bar, beta-bar, gamma-bar; m(U,V) for gamma(UxV).
%  - p. 22: pi(C,E), gamma(C,E) kept (shared with 161-4).
%  - pp. 24-30: nu_lambda for the number of ends (p. 26 writes a letter read
%    beta); g(alpha) for the genus the page calls gamma(alpha); Mod(S) for the
%    page's G = Isot_o(S,S), « groupe de Teichmüller »; Pi = pi_1(S_rho) for the
%    page's G_rho on p. 27 (which also calls Isot_o(S_rho,S_rho) G_rho).
%  - pp. 32-33: H = G_a (arc stabiliser), H' = G_s (vertex stabiliser), as p. 32.
%  Letters S, G, Gamma, alpha, pi mean different things in different parts;
%  the spine section lists the collisions.
%
% Substantive departures, with pages.
%  p. 3   O(E) « doublement transitive » on (u,v): transitive, stabilisers of order 2.
%  p. 4   existence of alpha needs theta > pi/3 (page: >=; theta = pi/3 gives alpha = 0).
%  p. 4   « rayon 4 sin^2(alpha/2) »: that is the squared chord; radius l = 2 sin(alpha/2).
%  p. 5   (4) used twice; renamed l, l'. G_C for G'_C twice (primes supplied).
%  p. 5   Lemma « de sommet u »: the summit is x_0.
%  p. 5   gamma^u_{theta/2} named, rotation of angle -theta/2 meant.
%  p. 6   computation stops; the six relations are completed here (ours) and
%         reduced to eta^2 + eta - 1 = 0 (the link to pp. 12, 14 is ours).
%  p. 7   « I une arête »: d an edge through x.
%  pp. 8-9  Aut(vertex-edge graph) and Aut(face-edge graph) are Gamma (order
%         120), not Gamma' as the page writes; his own reduction (fix u, x_0,
%         x_1) and p. 10 both require Gamma. Checked by brute force.
%  p. 13  c) « éléments de prod' » -> of its image under m.
%  p. 15  0_I -> 0_J, 0_K -> 0_L, 1_K -> 1_L.
%  (OM)   not defined in the folder; read as frames from what the proofs use
%         (infinite distributivity in L used silently).
%  p. 17  « 4pi/5 < 2pi/3 » false. p. 18 cos alpha chain wrong; p. 19 boxed
%         cos alpha is cos(alpha/2). All footnoted, not used.
%  p. 18  injectivity of G -> S_5 completed (ours, class sizes).
%  p. 21  margin c^2 + c - 1 = 0 for c = cos'(D_i,D_j): false (c = -6/5); the
%         golden relation is eta's. First equation for a, b: sign of b gamma^2.
%         Base: 2 and 5 invertible supplied; Z[1/4] kept in a footnote.
%  p. 22  p_C a functor, not a morphism of topoi. The Question: yes, classical.
%  p. 25  margin « b_1 + 1 components »: true iff g(alpha) = 0; general
%         formula |I(alpha)| = b_0 + b_1 - 2 g(alpha) (ours). Prop bounds
%         (S^1 x [1,2[, {1}, ]0,2[) inconsistent as read; stated correctly.
%  p. 26  genus « sum g_lambda + gamma(alpha) — vérifier »: correct formula adds
%         b_1 of the gluing graph (ours).
%  p. 27  Isot a torsor under Autext: under Mod^+, which injects into Out with
%         image of index 2 (genus >= 1, DNB). p. 28 « est un isom. » idem.
%  p. 28  normal case « nu_lambda nuls »: nu_lambda <= 1 (0 when S compact).
%         H = 1 in the normal case: last step ours (DNB injectivity).
%  p. 28  exponent of S^2 read 2 against the hand (1).
%  p. 32  G « ≃ » Perm(S) x Perm(A): a subgroup.
%  p. 33  labels G_a, G_s under H', H swapped; condition sigma-dot not in H'/H
%         (edges not loops) supplied.
%
% Announced and never established on the pages: p. 21 a)-c); p. 22 the
% Question (answered by the literature, and also worked in 161-4); p. 29 the
% Question (left open here); p. 30 exactness (« ?? »); p. 13 e) and the
% infinite version's conclusion; pp. 12-15 the binary hypotheses a)-c) are
% on no page of this folder.
\documentclass[11pt,a4paper]{article}
% Le préambule commun à toutes les transcriptions du fonds Grothendieck.
%
% Il tient en une idée : **l'incertitude se compose, elle ne se résout pas.**
% Une transcription qui lisse un mot douteux a détruit la seule chose pour
% laquelle on la faisait — un lecteur doit pouvoir savoir, à chaque endroit,
% ce qui était sur le feuillet et ce qui a été ajouté. Les six macros qui
% suivent sont donc l'appareil critique, et non de la décoration.
%
% Les mêmes commandes sont reconnues par `scripts/render.mjs`, qui produit la
% vue de lecture du site. Une macro ajoutée ici doit l'être aussi là-bas,
% faute de quoi le rendu échouera bruyamment — ce qui est voulu : un rendu
% silencieusement faux est pire qu'un rendu absent.

\ProvidesPackage{grothendieck}[2026/08/08 Transcriptions du fonds Grothendieck]

\usepackage{amsmath,amssymb,amsthm}
\usepackage{mathrsfs}
\usepackage{stmaryrd}
% Les diagrammes commutatifs sont partout dans ce fonds : c'est la langue
% même de la géométrie algébrique de Grothendieck, et une transcription qui
% les rendrait en prose ne transcrirait rien.
\usepackage{tikz-cd}
% Français par défaut — c'est la langue des feuillets — mais l'anglais est
% chargé aussi : la traduction et le résumé sont des documents anglais, et la
% typographie française (espaces avant les deux-points) y serait une faute.
% Ces documents basculent par \selectlanguage{english} en tête de corps.
\usepackage[english,french]{babel}
\usepackage{fontspec}
\usepackage{geometry}
\geometry{a4paper,margin=25mm,marginparwidth=28mm}
\usepackage{marginnote}
\usepackage{xcolor}
% ulem, not soul. `soul`'s \st and \ul are built on a character-by-character
% scan that XeLaTeX's font machinery breaks: the document fails with an
% undefined control sequence rather than degrading. ulem does the same two jobs
% and is Unicode-engine safe. `normalem` stops it hijacking \emph, which the
% transcriptions use for Grothendieck's own emphasis.
\usepackage[normalem]{ulem}
\usepackage{hyperref}
\hypersetup{colorlinks=true,linkcolor=black,urlcolor=black,pdfborder={0 0 0}}

% --- Métadonnées du lot -----------------------------------------------------
% Reprises telles quelles par le script de rendu, qui les affiche en tête de
% la vue de lecture. Elles fixent ce que le fichier prétend couvrir : sans
% elles, une transcription flotte sans point d'attache dans un fonds de seize
% mille feuillets.

\newcommand{\folder}[1]{\gdef\grothFolder{#1}}
\newcommand{\batch}[1]{\gdef\grothBatch{#1}}
\newcommand{\leaves}[2]{\gdef\grothFirst{#1}\gdef\grothLast{#2}}
\newcommand{\foldertitle}[1]{\gdef\grothFoldertitle{#1}}
\gdef\grothFolder{?}\gdef\grothBatch{?}\gdef\grothFirst{?}\gdef\grothLast{?}\gdef\grothFoldertitle{}

% --- Le résumé --------------------------------------------------------------
%
% Il ouvre la lecture modernisée, sous le titre « Résumé », et il est composé
% en petit. Ce n'est pas un ornement : c'est le seul passage du document qui
% ne soit pas des mathématiques, et un lecteur doit pouvoir le reconnaître
% avant de l'avoir lu — puis le sauter s'il connaît déjà le sujet, ou s'y
% arrêter s'il ne le connaît pas. Le filet qui le suit marque la reprise du
% corps.

\newenvironment{resume}
  {\small\setlength{\parskip}{0.45em}}
  {\par\medskip\hrule\medskip}

% --- Mots-clés ---------------------------------------------------------------
%
% En anglais, en clôture du résumé de la lecture modernisée : le vocabulaire
% moderne sous lequel chercher ce que les feuillets construisent. Le manifeste
% du site (`npm run manifest`) les extrait pour étiqueter la cote entière —
% cette ligne est la source unique des tags, il n'y a pas de fichier à côté.
\newcommand{\keywords}[1]{\par\smallskip\noindent
  {\footnotesize\textsc{Keywords} --- \textit{#1}}\par}

% --- L'appareil critique ----------------------------------------------------

% Le numéro de feuillet, en marge. C'est l'ancre qui permet de régler un
% désaccord sur une formule en regardant *un* feuillet plutôt qu'une chemise
% de six cents. Le site s'en sert aussi pour tourner les pages du fac-similé
% quand on fait défiler la transcription.
\newcommand{\leaf}[1]{%
  \marginnote{\footnotesize\color{gray!70}\textsf{#1}}%
  \label{leaf:#1}\ignorespaces}

% Illisible : on ne devine pas. Un mot inventé qui se lit comme les autres est
% le pire résultat possible.
\newcommand{\ill}{\textcolor{red!60!black}{[\,\ldots\,]}}

% Lecture douteuse : le mot est donné, et signalé comme incertain.
\newcommand{\uncertain}[1]{\uline{#1}}

% Ajout de l'éditeur — une lettre manquante, un mot sous-entendu.
\newcommand{\add}[1]{\textcolor{blue!50!black}{[#1]}}

% Biffé par Grothendieck lui-même. Ce qu'il a rayé fait partie du document :
% une rature dit souvent où la pensée a changé de direction.
\newcommand{\struck}[1]{\sout{#1}}

% Note du transcripteur, sur l'état matériel du feuillet.
\newcommand{\note}[1]{\par\smallskip{\small\color{gray!60!black}\textit{[#1]}}\par\smallskip}

% Note marginale de Grothendieck, distincte du corps du texte.
\newcommand{\marginal}[1]{\par\smallskip\noindent\hspace{1em}%
  {\small\color{gray!60!black}#1}\par\smallskip}

% --- Titre ------------------------------------------------------------------

\AtBeginDocument{%
  \begin{center}
    {\large\bfseries Fonds Alexandre Grothendieck --- cote n\textsuperscript{o}~\grothFolder}\\[2pt]
    {\itshape\grothFoldertitle}\\[4pt]
    {\small feuillets \grothFirst--\grothLast\ (lot \grothBatch)}\\[2pt]
    {\footnotesize\color{gray!60!black}Transcription établie d'après le fac-similé numérisé
     par l'Université de Montpellier.\\
     Les lectures incertaines sont soulignées, les illisibles notées [\,\ldots\,],
     les ajouts entre crochets.}
  \end{center}
  \vspace{1em}}

\endinput


\folder{161-6}
\pages{1}{33}
\dating{[à partir de 1973-vers 1977] — le groupe « Dossiers rassemblés par Grothendieck sur différentes thématiques » (161-1 à 162-6) est daté [après 1961-vers 1977]}
\watermark{Édition de démonstration\\interprétation personnelle de l'œuvre}
\foldertitle{Graphes, icosaèdre [etc.] : notes manuscrites (s.d.) — lecture modernisée du dossier entier}

\begin{document}

\section*{Résumé}

\begin{resume}
La chemise porte deux mots, « Graphes, icosaèdre », et l'inventaire ajoute
prudemment « etc. ». Les feuillets qu'elle contient ne forment pas un texte
suivi : une mise au net à l'encre, des brouillons, quatre pages d'un carnet à
spirale, quelques feuilles volantes. Ils traitent de trois sujets, et d'un
quatrième en passant.

Le premier est l'un des plus anciens objets des mathématiques, l'icosaèdre
régulier, le solide à vingt faces triangulaires. Une mise au net soignée le
construit à partir de presque rien : deux vecteurs unitaires orthogonaux, un
angle de 72 degrés, un peu de trigonométrie. On obtient douze points sur la
sphère, et l'on montre que les rotations qui les préservent sont au nombre de
soixante, chacune déterminée par l'image d'une seule arête orientée. Puis
vient le mouvement qui fait l'intérêt de ces pages. Le solide construit, la
géométrie est mise de côté : tout ce qu'on a démontré sur ce groupe se lit sur
le seul schéma d'incidence --- quels sommets sont reliés à quels sommets. Le
solide devient un \emph{graphe}. Une chose pourtant échappe au graphe :
l'orientation de l'espace, ce qui distingue une rotation de son reflet dans un
miroir. Pour la retrouver, il faut ajouter en chaque sommet un ordre
circulaire sur les arêtes qui en partent, comme les heures sur un cadran. La
page s'arrête sur ce souhait.

Autour de cette mise au net gravitent des calculs sur le nombre
$2\cos(2\pi/5)$, racine de $t^{2} + t - 1 = 0$ --- l'inverse du nombre d'or ---,
faits dans l'anneau engendré par les racines cinquièmes de l'unité, et une
page qui transporte l'icosaèdre en géométrie algébrique : un icosaèdre non
plus dans l'espace euclidien mais relativement à une forme quadratique sur une
base quelconque, dont elle annonce, sans le démontrer, que le groupe de
symétries reste en chaque point le groupe alterné à cinq lettres.

Le deuxième sujet prend l'ordre circulaire aux sommets pour point de départ,
dans une tout autre généralité : un graphe quelconque dessiné sans croisement
sur une surface. Posez un graphe en fil de fer sur une table et collez un
ruban de papier le long de chaque fil, en recollant les rubans aux carrefours
dans l'ordre où les fils y arrivent : vous obtenez une surface percée de
trous. Grothendieck énonce que cette surface-ruban ne dépend que des ordres
circulaires, et montre que classer les dessins revient à classer les façons de la
compléter en une surface, puis à comprendre comment les déformations de la
surface agissent sur eux. Il y retrouve le groupe des classes
d'homéomorphismes d'une surface, qu'il appelle groupe de Teichmüller, et
laisse une question ouverte. Deux pages enfin décrivent les graphes dont le
groupe de symétries agit transitivement sur les sommets et les arêtes : trois
sous-groupes emboîtés et un élément d'ordre deux suffisent à les coder, et
l'icosaèdre en est l'exemple --- rapprochement qui est le nôtre.

Le sujet traité en passant est le produit. Une suite au crayon démontre qu'un
certain ensemble ordonné est la « somme » de deux autres, et ce qu'elle
manipule est ce qu'on appelle aujourd'hui un \emph{frame} : l'ensemble
ordonné des ouverts d'un espace, pris pour lui-même. Une page barrée pose, pour
les topos, la question voisine du produit d'un topos par un topos de
préfaisceaux --- question que le dossier 161-4 travaille aussi, dans les mêmes
notations.

Ce que ces pages montrent de sa manière est net sur l'icosaèdre : une
construction faite avec des sinus et des cosinus est aussitôt relue comme une
structure combinatoire, et la question devient de savoir ce que la
combinatoire voit et ce qu'elle ne voit pas. La datation de l'inventaire,
« à partir de 1973 », place ces feuillets à Montpellier. Les noms sous
lesquels chercher la suite sont ceux de graphe enrubanné (\emph{ribbon
graph}), de système de rotation, de groupe modulaire d'une surface, de graphe
arc-transitif, de frame et de locale. Les graphes sur les surfaces et
l'icosaèdre sont aussi des objets qu'étudie l'\emph{Esquisse d'un programme}
(1984), sous le nom de « cartes » et de « dessins d'enfants » ; ces pages ne la
citent pas.

\keywords{regular icosahedron, icosahedral group, golden ratio, real cyclotomic field, equiangular lines, frame, coproduct of frames, product of topoi, presheaf topos, ribbon graph, rotation system, mapping class group, Dehn–Nielsen–Baer theorem, arc-transitive graph, coset graph}
\end{resume}

\pagerange{1}{33}
\section*{Le fil du dossier, et les conventions}

\emph{Les pièces.} Le dossier est fait de feuillets de nature différente,
sans pagination de sa main. Les pages 2, 11, 20, 23 et 31 sont blanches.

\begin{enumerate}
  \item \emph{Page 1}, la chemise. Sur l'onglet, au crayon : « Graphes,
  icosaèdre ». En tête, un « 6 » cerclé, un mot qui se lit peut-être
  « Groupes », et entre parenthèses une mention de prêt : « Notes sur formes
  paradoxiques prêtées à Labarde », ou « à Laborde » : les deux mots décisifs
  sont d'une lecture incertaine, et dans le nom la main ne distingue pas le
  a du o.\footnote{Le dossier 162-1 du même fonds est un registre de prêts.
  Rien ne permet de rattacher cette mention à ce registre : la
  transcription de 162-1, parcourue pour ce travail, ne porte le nom sous
  aucune de ses lectures, et les « formes paradoxiques » ne correspondent à
  aucune de ses entrées. On ne suppose aucun lien.}
  \item \emph{Pages 3 à 10}, encre sur papier réglé, titrée « L'icosaèdre » :
  une mise au net (section I).
  \item \emph{Pages 14 et 12}, encre : calculs sur les racines cinquièmes de
  l'unité, à lire dans cet ordre (section III).
  \item \emph{Pages 15, 13 et pied de la page 12}, crayon : une démonstration
  sur des ensembles ordonnés (section V). Les pages 15 et 13 sont barrées
  chacune d'un long trait oblique, les deux lignes du pied de la page 12 ne le
  sont pas ; la page 13 et ces deux lignes sont écrites tête-bêche par rapport
  à l'encre, la page 15 dans le même sens qu'elle.
  \item \emph{Pages 16 à 19}, quatre feuillets d'un carnet à spirale :
  croquis, le calcul des pages 3 et 4 avec d'autres lettres, le groupe $\mathfrak{A}_5$
  (section II).
  \item \emph{Page 21}, l'icosaèdre relativement à une forme quadratique
  (section IV).
  \item \emph{Page 22}, écrite dans l'autre sens et barrée d'une croix : le
  produit d'un topos par un topos de préfaisceaux (section VI).
  \item \emph{Pages 24 à 30} : plongements d'un graphe dans une surface
  orientée (section VII).
  \item \emph{Pages 32 et 33} : graphes à groupe transitif (section VIII).
\end{enumerate}

\emph{Le fil.} Trois sujets, que rien sur les pages ne relie explicitement.
L'icosaèdre (I à IV) est construit, puis réduit à son graphe, et la page 10
demande comment retrouver combinatoirement les rotations parmi les
symétries ; la réponse qu'elle esquisse --- des ordres circulaires
« cohérents » aux sommets --- est exactement la donnée $\omega(\varphi)$
dont partent les pages 24 à 30, et l'icosaèdre est un cas des graphes des
pages 32 et 33. Ces deux rapprochements sont les nôtres ; ils expliquent
peut-être pourquoi ces feuillets sont réunis sous « Graphes, icosaèdre »,
mais ils ne sont écrits nulle part. Le produit (V et VI) est un sujet à part :
la suite au crayon établit une propriété universelle de somme dans une
catégorie d'ensembles ordonnés, la page 22 pose la question du produit de
deux topos ; que la seconde soit la version topos de la première est encore
notre lecture.

\emph{Notations.} Chaque partie garde les lettres de ses pages, sauf là où
une même lettre y désigne deux objets. Les renommages sont les suivants,
chacun signalé en note au point où il se produit : la rotation que les pages
5 et 6 appellent $\gamma$ est notée $r$, et $R_{a,t}$ désigne la rotation
d'angle $t$ autour du vecteur unitaire $a$ ; le nombre
$\zeta + \zeta^{-1} = 2\cos(2\pi/5)$, que les pages appellent tour à tour
$\gamma$, $\eta$, $t$ ou $c$, est noté $\eta$ ; dans la suite au crayon,
l'application que la page appelle $\gamma$ est notée $m$ ; aux pages 24 à 30,
le genre que la page appelle $\gamma(\alpha)$ est noté $g(\alpha)$, et le
groupe que la page appelle $G$ est noté $\mathrm{Mod}(S)$. Restent des
collisions entre parties, qu'on ne corrige pas mais qu'il faut avoir en
tête : $S$ est la sphère unité en I, un schéma de base en IV, une surface en
VII, l'ensemble des sommets en VIII ; $G$ est $\mathrm{O}(E)$ en I, un groupe
abstrait en II et VIII ; $\alpha$ est un angle en I, un système d'ordres
circulaires en VII ; $\pi$ est une partition en VII.

\emph{Ce que le dossier annonce sans l'établir.} Les énoncés a), b), c) de la
page 21 ; la réponse à la Question de la page 22 ; la Question de la page 29 ;
l'exactitude de la suite de la page 30, marquée « ?? » ; la conclusion de la
version infinie de la page 13. Les hypothèses de la suite au crayon, dans sa
forme à deux facteurs, sont sur une page qui n'est pas dans ce dossier.

\pagerange{3}{10}
\section*{I. L'icosaèdre (pages 3 à 10)}

\pagerange{3}{4}
\subsection*{Le triangle équilatéral d'angle $\theta$ (pages 3 et 4)}

Soit $E$ un espace euclidien orienté de dimension $3$, de produit scalaire
$x \cdot y$, et $S$ sa sphère unité.\footnote{En marge de la page 3 : « espace
vectoriel muni d'une forme quadratique non dégénérée, qu'on supposera définie
positive au besoin ». C'est le cadre de la page 21, où l'icosaèdre est défini
relativement à une forme quadratique quelconque.} Pour $u \in S$, on note
$\Sigma_u = u^{\perp} \cap S$ l'ensemble des directions tangentes unitaires en
$u$ ; pour $v \in \Sigma_u$, on pose $w = u \wedge v$, de sorte que
$(u, v, w)$ est une base orthonormée directe. Le groupe $\mathrm{SO}(E)$ agit
simplement transitivement sur l'ensemble des couples $(u, v)$ avec
$u \in S$, $v \in \Sigma_u$ ; le groupe
$\mathrm{O}(E) = \mathrm{SO}(E) \times \{\pm 1\}$ y agit transitivement, le
stabilisateur d'un couple étant d'ordre $2$, engendré par la symétrie par
rapport au plan de $u$ et $v$.\footnote{La page dit que $\mathrm{O}(E)$ opère
« de façon doublement transitive » ; ce qu'elle veut dire est ce qui est
écrit ici --- la double transitivité, au sens usuel, est autre chose. Elle
écrit aussi « $(u,v) \in S$ » là où il faut $S \times S$.}

Le grand cercle issu de $u$ dans la direction $v$ est paramétré par
$x(u,v;\alpha) = u\cos\alpha + v\sin\alpha$, et les directions tangentes en
$u$ s'écrivent $v' = v\cos\theta + w\sin\theta$. Pour $\theta$ fixé, le
triangle de sommets $u$, $x = x(u,v;\alpha)$, $x' = x(u,v';\alpha)$ est
isocèle, et
\[
\|x - u\|^{2} = \|x' - u\|^{2} = 4\sin^{2}\tfrac{\alpha}{2},
\qquad
\|x - x'\|^{2} = \sin^{2}\alpha\,\|v - v'\|^{2}
= 16\sin^{2}\tfrac{\alpha}{2}\cos^{2}\tfrac{\alpha}{2}\sin^{2}\tfrac{\theta}{2}.
\]
Pour $0 < \alpha < \pi$, il est donc équilatéral si et seulement si
\[
(1)\quad \cos\frac{\alpha}{2} = \frac{1}{2\sin(\theta/2)},
\qquad\text{soit}\qquad
(1')\quad \cos\alpha = \frac{\cos\theta}{1 - \cos\theta}.
\]
Pour $0 < \theta \leqslant \pi$, un tel $\alpha$ existe si et seulement si
$\theta > \pi/3$.\footnote{La marge de la page 4 écrit
$|\sin\frac{\theta}{2}| \geqslant \frac12$ et conclut
$\theta \geqslant \frac{\pi}{3}$ ; pour $\theta = \pi/3$ on trouve $\alpha = 0$
et le triangle dégénère en un point, d'où l'inégalité stricte. La page écrit
aussi $\cos\frac{\alpha}{2} = \pm\frac{1}{2\sin\theta/2}$ ; le signe $-$ donne
$2\pi - \alpha$, c'est-à-dire le même cercle parcouru dans l'autre sens, et
l'ajout interlinéaire « choisissant la détermination principale de $\alpha$ »
l'écarte.} Les trois angles sphériques du triangle sont alors égaux à
$\theta$, et c'est le seul triangle équilatéral inscrit dans $S$, d'angle
$\theta$, ayant un sommet en $u$, un deuxième dans la direction $v$, et tel
que $(u, x, x')$ soit direct.

On fixe désormais
\[
(2)\quad \theta = \frac{2\pi}{5},
\qquad\text{d'où}\qquad
\cos\alpha = \frac{1}{\sqrt{5}},\quad \sin\alpha = \frac{2}{\sqrt{5}},\quad \cos 2\alpha = -\frac{3}{5},
\]
les deux dernières valeurs étant notées en marge de la page 4. On abrège
$c = \cos\alpha$ et $s = \sin\alpha$.

\pagerange{4}{5}
\subsection*{Les douze points, et les calottes (pages 4 et 5)}

Pour $u$, $v$ fixés, on pose, pour $i \in \mathbb{Z}/5\mathbb{Z}$,
\[
v_i = v\cos i\theta + w\sin i\theta,\qquad
x_i = u\cos\alpha + v_i\sin\alpha,\qquad
x'_i = -x_i,\qquad u' = -u .
\]
Les six points $u, x_0, \ldots, x_4$ sont distincts --- leurs projections sur
$u^{\perp}$ sont $0$ et les $s\,v_i$ --- et ont une coordonnée positive le
long de $u$, puisque $c > 0$ ; les douze points $u, u', x_i, x'_i$ sont donc
distincts. Leur ensemble est l'\emph{icosaèdre} $I = I_{u,v}$ défini par
$(u,v)$. Les $x_i$ sont en ordre circulaire sur le cercle de $S$ de centre $u$
et de rayon (rectiligne) $\ell = 2\sin\frac{\alpha}{2}$.\footnote{La page 4
donne pour rayon $4\sin^{2}\frac{\alpha}{2}$, qui est le carré de la distance
calculée page 3.}

On note
\[
\ell = 2\sin\frac{\alpha}{2} = \|a - b_i\|,\qquad \ell' = 2\cos\frac{\alpha}{2} = \|a + b_i\|,
\]
les deux distances qui interviennent entre le centre $a$ d'une calotte et ses
points périphériques $b_i$ ou leurs opposés ; $\ell^{2}$ et $\ell'^{2}$ sont
les racines de $\lambda^{2} - 4\lambda + 4\sin^{2}\alpha = 0$.\footnote{La page
numérote (4) la formule de $\ell$, alors que (4) numérote déjà, une page plus
haut, la définition des $x_i$ ; (4') est la formule de $\ell'$. On nomme les
deux longueurs sans les numéroter.}

Un \emph{triangle icosaédral} est un ensemble de trois points de $S$ deux à
deux à distance $\ell$ ; de façon équivalente, l'image de $\{u, x_0, x_1\}$
par une rotation. Une \emph{calotte d'icosaèdre} est l'image de
$\{u, x_0, \ldots, x_4\}$ par un élément de $\mathrm{O}(E)$. Il revient au
même de demander un ensemble de six points qu'on puisse écrire
$(a; b_0, \ldots, b_4)$ de sorte que les triangles $(a, b_i, b_{i+1})$,
$0 \leqslant i \leqslant 3$, soient icosaédraux : la rotation d'angle $\pm\theta$
autour de $a$ envoie alors chaque $b_i$ sur $b_{i+1}$, et, parce que
$5\theta = 2\pi$, le triangle $(a, b_4, b_0)$ est lui aussi icosaédral. Le
point $a$, le seul qui soit à distance $\ell$ des cinq autres, est le
\emph{sommet} de la calotte ; les $b_i$ en sont les \emph{points
périphériques}. L'orientation de $E$ munit ceux-ci d'un ordre circulaire
canonique --- celui dans lequel les rencontre une rotation directe autour de
$a$ ---, qui en fait un torseur sous $\mathbb{Z}/5\mathbb{Z}$.

Le stabilisateur $G_C$ d'une calotte $C$ dans $G = \mathrm{O}(E)$ est un
groupe diédral d'ordre $10$. Son intersection $G'_C$ avec
$G' = \mathrm{SO}(E)$ est le groupe des rotations autour de $a$ d'angle
multiple de $\theta$, canoniquement isomorphe à $\mathbb{Z}/5\mathbb{Z}$ grâce
à l'orientation, et il agit simplement transitivement sur les points
périphériques, ce qui redonne leur ordre circulaire. Les cinq autres éléments
sont d'ordre $2$ : ce sont les symétries $g_i$ par rapport aux plans
$\langle a, b_i \rangle$, et $\{1, g_i\}$ est le stabilisateur de $b_i$ dans
$G_C$. Ainsi $G_C = G'_C \rtimes \mathbb{Z}/2\mathbb{Z}$, le facteur
$\mathbb{Z}/2\mathbb{Z}$ agissant sur $G'_C$ par
inversion.\footnote{Aux deux endroits où la page écrit « le sous-groupe
$G_C$ » et « opérant sur $G_C$ », le sens demande $G'_C$ ; on ne voit pas le
prime sur la page. La page dit « les cinq points de
$G_C \cap \complement G'_C$ » pour les cinq éléments, et la phrase qui les
décrit reste inachevée ; « produit ½ direct » est le produit semi-direct.}

Jusqu'ici, la page le note, $\theta = 2\pi/5$ n'a joué aucun rôle : on aurait
pu prendre $\theta = 2\pi/m$. Ce choix intervient avec le lemme suivant.

\pagerange{5}{7}
\subsection*{Le lemme de la calotte voisine (pages 5 à 7)}

\textbf{Lemme.} L'ensemble $\{x_0, x_1, u, x_{-1}, x'_2, x'_{-2}\}$ est une
calotte de sommet $x_0$, et l'ordre circulaire canonique de ses points
périphériques est $(x_1, u, x_{-1}, x'_2, x'_{-2})$.\footnote{La page écrit la
calotte $(x_0; x_1, u, x_{-1}, x'_2, x'_{-2})$, sommet en tête, puis « calotte
d'icosaèdre de sommet $u$ » ; la lettre après « sommet » est un lapsus pour
$x_0$. On a vérifié que l'ordre écrit est bien l'ordre direct autour de
$x_0$.}

La page ébauche d'abord une démonstration par les distances : la symétrie par
rapport au plan $\langle u, x_0 \rangle$ échange $x_1$ et $x_{-1}$, $x'_2$ et
$x'_{-2}$, et ramène le lemme, à l'orientation près, à vérifier que
$x_0 - x'_2$, $x_{-1} - x'_2$ et $x'_2 - x'_{-2}$ sont de longueur $\ell$ ---
la dernière étant évidente. Elle préfère ensuite une voie « plus astucieuse ».
Soit
\[
r = R_{w,\alpha} \circ R_{u,-\theta/2},
\]
composé de la rotation d'angle $-\theta/2$ autour de $u$ et de la rotation
d'angle $\alpha$ autour de $w$.\footnote{La page appelle cette rotation
$\gamma = \gamma^{w}_{\alpha}\gamma^{u}_{-\theta/2}$, après avoir nommé
$\gamma^{u}_{\theta/2}$, avec l'indice sans signe, « la rotation d'angle
$-\theta/2$ autour de $u$ ». On la note $r$ : la lettre $\gamma$ désigne dans
ce dossier trois autres objets.} Sa matrice dans la base $(u, v, w)$ est
\[
r = \begin{pmatrix}
\cos\alpha & -\sin\alpha\cos\frac{\theta}{2} & -\sin\alpha\sin\frac{\theta}{2} \\
\sin\alpha & \cos\alpha\cos\frac{\theta}{2} & \cos\alpha\sin\frac{\theta}{2} \\
0 & -\sin\frac{\theta}{2} & \cos\frac{\theta}{2}
\end{pmatrix},
\]
et le lemme résulte des six relations
\[
r u = x_0,\quad r x_0 = x'_2,\quad r x_1 = x'_{-2},\quad r x_2 = x_1,\quad r x_3 = u,\quad r x_4 = x_{-1},
\]
puisque $r$ est une rotation, qui transporte la calotte $(u; x_0, \ldots, x_4)$
et son ordre canonique.

La page vérifie la première, écrit $r x_0$ et pose les trois égalités à
prouver pour la deuxième, commence $r x_1$ et s'arrête. On achève le
calcul.\footnote{La vérification qui suit est la nôtre ; la page 6 s'arrête
sur « $\gamma x_1 = \gamma u \cos\alpha + \ldots = $ », le reste de la page
étant blanc. Les trois égalités qu'elle pose pour $r x_0 = x'_2$ sont
exactes : ce sont, dans l'ordre, les coordonnées selon $u$, $v$ et $w$
ci-dessous, pour $i = 0$.} La rotation $R_{u,-\theta/2}$ envoie $v_i$ sur le
vecteur unitaire d'angle $\beta_i = (2i - 1)\pi/5$ dans le plan $(v, w)$, d'où
\[
r x_i = \bigl(c^{2} - s^{2}\cos\beta_i\bigr)\,u + cs\,(1 + \cos\beta_i)\,v + s\sin\beta_i\,w .
\]
On compare aux coordonnées de $x_j = c\,u + s\cos j\theta\,v + s\sin j\theta\,w$
et de $x'_j = -x_j$. Les coordonnées selon $w$ concordent d'elles-mêmes ; pour
$i = 3$, $\beta_3 = \pi$ et $r x_3 = u$ ; pour les quatre autres valeurs de
$i$, en utilisant $s^{2} = 1 - c^{2}$, les deux autres coordonnées se ramènent
aux deux identités
\[
\text{(A)}\quad c\,(1 - \cos\tfrac{2\pi}{5}) = \cos\tfrac{2\pi}{5},
\qquad
\text{(B)}\quad c\,(1 + \cos\tfrac{\pi}{5}) = \cos\tfrac{\pi}{5},
\]
(A) pour $i = 2$ et $i = 4$, (B) pour $i = 0$ et $i = 1$. L'identité (A) est
(1'). Posons $\eta = 2\cos\frac{2\pi}{5}$ et $p = 2\cos\frac{\pi}{5}$ ; alors
$p = \eta + 1$ et $p\,\eta = 1$, deux formes de la relation
$\eta^{2} + \eta - 1 = 0$, et, compte tenu de (A), l'identité (B) équivaut à
$\eta + p\eta = p$, qui en découle. C'est là, et seulement là, que
$\theta = 2\pi/5$ intervient : par l'équation du nombre d'or, que les pages 14
et 12 établissent pour elles-mêmes, sans rapport écrit avec
celle-ci.\footnote{Le rapprochement entre la vérification de la page 6 et
les calculs des pages 12 et 14 est le nôtre. Ces derniers sont écrits dans
un autre contexte (section III), mais ils calculent exactement ce qui manque
ici.}

\textbf{Corollaire.} Toute $g \in \mathrm{O}(E)$ qui envoie la calotte
$(u; x_0, \ldots, x_4)$ sur celle du lemme préserve l'icosaèdre $I$. En
effet, pour l'une et l'autre calotte $C$, on a $I = C \cup (-C)$, et $g$ est
linéaire. En particulier $r$ préserve $I$.

\pagerange{7}{8}
\subsection*{Le groupe de l'icosaèdre (pages 7 et 8)}

Soit $\Gamma$ le stabilisateur de $I$ dans $\mathrm{O}(E)$, et
$\Gamma' = \Gamma \cap \mathrm{SO}(E)$. Comme $-1 \in \Gamma$ et
$-1 \notin \Gamma'$, on a $\Gamma = \Gamma' \times \{\pm 1\}$.

Les \emph{arêtes} de $I$ sont les trente paires
\[
\{u, x_i\},\ \{u', x'_i\},\ \{x_i, x_{i+1}\},\ \{x'_i, x'_{i+1}\},\ \{x_i, x'_{i-2}\},\ \{x'_i, x_{i-2}\}\qquad (i \in \mathbb{Z}/5\mathbb{Z}),
\]
c'est-à-dire les paires de points distincts de $I$ à distance minimale
$\ell$. Ses \emph{faces} sont les vingt triples
\[
\{u, x_i, x_{i+1}\},\ \{u', x'_i, x'_{i+1}\},\ \{x_i, x'_{i-2}, x_{i+1}\},\ \{x'_i, x_{i-2}, x'_{i+1}\},
\]
c'est-à-dire les triples deux à deux reliés par une arête ; cinq d'entre eux
contiennent $u$. Tout point de $I$ est le sommet d'une unique calotte contenue
dans $I$, formée de lui-même et des cinq points qui lui sont reliés par une
arête.

\textbf{Théorème.} (i) $\Gamma'$ agit transitivement sur les douze sommets de
$I$. (ii) $\Gamma'$ agit simplement transitivement sur les soixante arêtes
orientées, c'est-à-dire sur les couples $(x, d)$ formés d'un sommet $x$ et
d'une arête $d$ issue de $x$. (iii) $\Gamma$ agit simplement transitivement
sur les cent vingt triples ordonnés $(x, y, z)$ qui forment une
face.\footnote{La page écrit « $(x,d)$ où $x$ est un sommet de l'icosaèdre,
$I$ une arête » : le $I$ est un lapsus pour « $d$ une arête issue de $x$ ».
Une note en marge, en partie illisible, ajoute qu'on peut dire la même chose
des triples d'arêtes formant une face.}

\emph{Démonstration.} La rotation $R_{u,\theta}$ préserve $I$ et permute
circulairement les $x_i$, et les $x'_i$ ; la rotation $r$ envoie $u$ sur
$x_0$ et $x_0$ sur $x'_2$, et, étant linéaire, $u'$ sur $x'_0$. Tous les
sommets sont donc congrus sous $\Gamma'$. Un élément de $\Gamma'$ qui fixe $u$
préserve l'ensemble des points de $I$ à distance $\ell$ de $u$ --- les autres
sont à distance $\ell' > \ell$ ou $2$ --- donc la calotte de sommet $u$ ;
inversement, une rotation qui préserve cette calotte $C$ préserve
$I = C \cup (-C)$. Le stabilisateur de $u$ dans $\Gamma'$ est donc $G'_C$,
simplement transitif sur les $x_i$, d'où (ii). Pour (iii), on se ramène par
(ii) aux triples $(u, x_0, z)$ ; il n'y a que deux faces contenant l'arête
$\{u, x_0\}$, de troisièmes sommets $x_1$ et $x_{-1}$, échangés par la
symétrie par rapport au plan $\langle u, v \rangle$, qui préserve $I$ ; et
une application linéaire qui fixe $u$, $x_0$, $x_1$, linéairement
indépendants, est l'identité.

\textbf{Corollaire.} $|\Gamma'| = 60$ et $|\Gamma| = 120$ : il y a trente
arêtes, donc soixante arêtes orientées ; et vingt faces, chacune avec six
ordres sur ses sommets.

En marge de la page 8, la liste des données équivalentes : un icosaèdre et
l'un de ses sommets ; un icosaèdre et une calotte contenue dans lui ; une
calotte ; un couple $(u, v)$ modulo les rotations de $v$ dans $u^{\perp}$
d'angle multiple de $\theta$. Autrement dit, l'espace de ces données est le
quotient de l'espace des couples $(u, v)$ --- un torseur sous
$\mathrm{SO}(E)$, qui est aussi un fibré principal sur $S$ de groupe
$\mathrm{SO}(2)$ --- par le sous-groupe $\mu_5$ du groupe structural :
c'est un fibré principal sur $S$ de groupe
$\mathrm{SO}(2)/\mu_5$.\footnote{La marge dit, entre crochets et d'une
lecture douteuse en plusieurs mots : « point du fibré principal sur $S$ de
groupe $\mathbb{T}$ quotient […] de celui formé des $(u,v)$ ». La
formulation ci-dessus est la nôtre.}

\pagerange{8}{10}
\subsection*{Deux graphes, et ce que la combinatoire ne voit pas (pages 8 à 10)}

Les sommets et les arêtes de $I$ forment un graphe $\mathcal{G}_1$ : douze
sommets, trente arêtes, cinq arêtes en chaque sommet. Le groupe $\Gamma$ y
agit fidèlement, puisque les sommets engendrent $E$.

\textbf{Proposition} (une « Remarque » de la page). Le groupe des
automorphismes de $\mathcal{G}_1$ est $\Gamma$, d'ordre $120$.\footnote{La page écrit : « Montrer que \emph{tout}
automorphisme du graphe est dans $\Gamma'$ », et, pour le second graphe,
« le groupe des automorphismes de ce graphe est encore $\Gamma'$ ». C'est
faux : $-1$ et les quinze symétries planes de $\Gamma$ sont des
automorphismes du graphe qui ne sont pas des rotations. La démonstration de
la page, elle, établit l'énoncé juste : elle se ramène aux automorphismes qui
fixent $u$, $x_0$ \emph{et} $x_1$, réduction qui repose sur la transitivité
de $\Gamma$ sur les faces ordonnées, point (iii), et non sur celle de
$\Gamma'$ ; la page 10 dit d'ailleurs « les actions de $\Gamma$ » et
cherche à caractériser $\Gamma'$ \emph{parmi} les symétries combinatoires,
ce qui n'aurait pas de sens si elles étaient toutes des rotations. Nous avons
vérifié par énumération que le graphe a exactement $120$ automorphismes.}

\emph{Démonstration.} Les faces sont les triangles du graphe, donc tout
automorphisme permute les triples ordonnés formant une face ; par (iii), il
suffit de montrer qu'un automorphisme qui fixe $u$, $x_0$, $x_1$ est
l'identité. Il fixe $x_2$, seul voisin commun de $u$ et $x_1$ autre que
$x_0$, et $x_{-1}$, seul voisin commun de $u$ et $x_0$ autre que $x_1$ ; puis
$x_3$, dernier voisin de $u$. La calotte de sommet $u$ est donc fixée point
par point ; de même celle de sommet $x_0$, puis, de proche en proche le long
des triangles $(u, x_i, x_{i+1})$, celles de sommet $x_i$. Cela donne tous les
sommets sauf $u'$, qui est fixé comme le seul restant. Le groupe des
automorphismes agit donc librement sur les $120$ faces ordonnées, ce qui
majore son ordre par $120 = |\Gamma|$.

Les faces et les arêtes forment un second graphe $\mathcal{G}_2$ : vingt
sommets, trente arêtes, trois arêtes en chaque sommet, une arête de $I$
reliant les deux faces qui la contiennent. C'est le graphe des sommets et
arêtes du dodécaèdre, dual de l'icosaèdre. Son groupe d'automorphismes est
encore $\Gamma$ : il suffit de voir qu'un automorphisme qui fixe une face,
disons $F_0 = \{u, x_0, x_1\}$, et chacune de ses trois voisines, est
l'identité. Il fixe alors $\{u, x_2, x_3\}$ et $\{u, x_3, x_4\}$, seul couple
de faces adjacentes entre elles et adjacentes respectivement à
$\{u, x_1, x_2\}$ et à $\{u, x_4, x_0\}$ ; donc chacune des voisines de $F_0$
a ses trois voisines fixées. De proche en proche, toutes les faces sont
fixées.\footnote{La page conclut : « par symétrie on voit que pour chacune
des faces adjacentes à $(u, x_0, x_1)$, ses faces et leurs arêtes sont
fixes. À partir de là la conclusion est triviale ». « Ses faces » est à lire
« ses faces adjacentes ».}

La page 10 tire la leçon. Les propriétés de l'action de $\Gamma$ sur les
sommets, les arêtes et les faces sont conséquences de la seule structure
combinatoire --- le graphe $\mathcal{G}_1$, le graphe $\mathcal{G}_2$, ou le
système des trois ensembles avec leurs incidences. Mais les rotations
$\Gamma'$, sous-groupe d'indice $2$, ne sont pas visibles sur le graphe seul.
La page propose de les caractériser comme les automorphismes qui respectent
les ordres circulaires des points périphériques des calottes, et note qu'on
peut définir combinatoirement \emph{deux} « systèmes cohérents » de tels
ordres, correspondant aux deux orientations de l'espace.

On peut préciser ce qu'est un système cohérent, dans les termes de la
section VII.\footnote{Ce paragraphe est le nôtre.} Un choix d'ordre
circulaire sur les cinq voisins de chaque sommet --- un \emph{système de
rotation} --- définit des cycles de bord, obtenus en tournant, à chaque
sommet atteint, vers l'arête suivante. Est cohérent un système dont les
cycles de bord sont les vingt triangles. Il y en a exactement deux,
échangés par tout élément de $\Gamma \smallsetminus \Gamma'$, et $\Gamma'$ est
le stabilisateur de l'un d'eux dans $\Gamma = \mathrm{Aut}(\mathcal{G}_1)$.
Qu'un graphe planaire $3$-connexe n'ait, à homéomorphisme près, qu'un
plongement dans la sphère est un théorème de Whitney (1933) ; rien n'indique
que Grothendieck le connaissait, et la page ne le cite pas. Un système
cohérent est exactement l'invariant $\omega(\varphi)$ que les pages 24 à 30
attachent au plongement de l'icosaèdre dans la sphère $S$.

\pagerange{16}{19}
\section*{II. Les feuillets du carnet (pages 16 à 19)}

Quatre feuillets d'un carnet à spirale, écrits en largeur et en travers,
couverts pour partie de croquis --- fleurs légendées « corolle », « racine »,
« tige », « cœur », « pétale », tétraèdres, un cube, une projection de
l'icosaèdre. On en retient trois choses.

\pagerange{17}{17}
\subsection*{Le calcul des pages 3 et 4, avec d'autres lettres (page 17)}

Avec d'autres lettres ($\rho_i$ pour $v_i$, $v'$ et $w'$ pour deux directions
tangentes faisant l'angle $2\pi/5$), la page 17 fait le calcul des côtés du
triangle et aboutit à la même condition (1), encadrée :
$\cos\frac{\alpha}{2} = 1/(2\sin\frac{\pi}{5})$. Elle note en marge que
$2\sin\frac{\pi}{5}$ est le côté du pentagone régulier inscrit dans le cercle
unité.\footnote{Parmi les inégalités éparses dans la marge des croquis, la
première, « $\frac{4\pi}{5} < \frac{2\pi}{3}$ », est fausse ; elles sont des
essais autour de la condition d'existence de la page 4 et ne servent nulle
part.}

\pagerange{18}{18}
\subsection*{Le groupe des rotations est $\mathfrak{A}_5$ (page 18)}

Soit $G$ un groupe d'ordre $60$ --- ici $\Gamma'$ --- agissant
transitivement sur un ensemble $I_{15}$ à quinze éléments, les paires
d'arêtes opposées de l'icosaèdre, c'est-à-dire les quinze axes passant par
les milieux des arêtes. Ces axes se groupent en cinq triples d'axes deux à
deux orthogonaux, et l'on peut munir chaque triple d'un ordre circulaire, de
façon invariante par $G$ : c'est l'automorphisme $\alpha$ d'ordre $3$ de
$I_{15}$ commutant à $G$ de la page, dont le quotient est un ensemble $J$ à
cinq éléments.\footnote{La page vérifie sur la figure --- une projection de
l'icosaèdre le long d'un axe, dont les sommets sont marqués $A_i$, $A'_i$ ---
que $\alpha^{3}(P) = P$ pour une paire $P$, en écrivant
$(AB, A'B') = \alpha(R) = \alpha^{2}(Q) = \alpha^{3}(P)$. Cette chaîne est
exacte : $\alpha(R) = \alpha(\alpha(Q)) = \alpha^{2}(Q)$, et
$\alpha^{3}(P) = P$ puisque $\alpha$ est d'ordre $3$.} Le groupe $G$ agit
transitivement sur $J$, avec des stabilisateurs d'ordre $12$ ; d'où un
morphisme $G \to \mathfrak{S}_5$. S'il est injectif, son image est un
sous-groupe d'indice $2$ de $\mathfrak{S}_5$, donc $\mathfrak{A}_5$, et
$G \simeq \mathfrak{A}_5$.

L'injectivité, la page commence à l'établir en passant par le normalisateur
d'un $2$-sous-groupe de Sylow, et s'arrête au bord du feuillet. On l'achève
ainsi.\footnote{Cet argument est le nôtre.} Le noyau $K$ est distingué dans
$G$ et contenu dans un stabilisateur, donc d'ordre au plus $12$ ; il est
réunion de classes de conjugaison. Or les classes non triviales de $\Gamma'$
ont $15$, $20$, $12$ et $12$ éléments --- les demi-tours autour des axes
d'arêtes, les rotations d'angle $\pm 2\pi/3$ autour des axes de faces, les
rotations d'angle $\pm 2\pi/5$ et celles d'angle $\pm 4\pi/5$ autour des axes
de sommets ---, ce qu'on lit sur les centralisateurs, d'ordres $4$, $3$, $5$
et $5$. Aucune ne tient dans $K \smallsetminus \{1\}$, qui a au plus onze
éléments ; donc $K = \{1\}$.

\pagerange{16}{19}
\subsection*{Deux lemmes sur les ensembles homogènes (pages 16 et 19)}

Au milieu des croquis, deux calculs qui sont exactement les outils des pages
32 et 33. Page 16 : pour un $G$-ensemble $E$ et un sous-groupe $H$,
\[
\mathrm{Hom}_G(G/H, E) \simeq E^{H},
\]
une application équivariante étant déterminée par l'image de la classe
neutre, qui doit être fixée par $H$. Page 19 : la multiplication à droite
par $a \in G$ passe au quotient $G/H$ si et seulement si $a^{-1}Ha \subset H$ ;
d'où, $G$ étant fini,
\[
\mathrm{Aut}_G(G/H) \simeq N_G(H)/H .
\]
La page 19 note aussi que le normalisateur d'un sous-groupe de Klein dans
$\mathfrak{A}_5$ est $(\mathbb{Z}/2\mathbb{Z})^{2} \rtimes \mathbb{Z}/3\mathbb{Z} \simeq \mathfrak{A}_4$,
le stabilisateur d'ordre $12$ ci-dessus.\footnote{La page 19 contient aussi
des essais trigonométriques qui ne s'accordent pas avec (1) : on y lit,
encadré, $\cos\alpha = 1/(2\sin\frac{2\pi}{10})$, qui est la valeur de
$\cos\frac{\alpha}{2}$, non de $\cos\alpha$. De même, page 18, la chaîne
$\cos\alpha = \ldots = \frac{1}{2\sin^{2}\theta} - 1 = 1 - 2\sin^{2}\theta = \cos 2\theta$
confond $\theta$ et $\theta/2$ et contient une égalité fausse ; la formule
juste est (1'). Ce sont des brouillons ; la mise au net a la formule juste.}

\pagerange{12}{14}
\section*{III. Le nombre d'or dans le corps des racines cinquièmes (pages 14 et 12)}

Les deux pages d'encre se lisent dans l'ordre 14, puis 12 : le calcul de la
norme de $\eta$, qui s'arrête page 14 sur un signe $=$, se termine en tête de
la page 12, qui en tire l'équation de $\eta$.\footnote{Cet ordre de lecture
est le nôtre ; la transcription ne le relève pas. Il est confirmé par la
suite au crayon des mêmes feuilles, qui se lit elle aussi 15 avant 13
(section V) : les deux feuilles semblent avoir été rangées dans l'ordre
inverse de celui où elles ont été écrites.}

\pagerange{14}{14}
\subsection*{Le quotient par l'inversion (page 14)}

Soit $A = \mathbb{Z}[\frac15]$ et $\mu_5^{*} \subset \mu_5$ le schéma des
racines primitives cinquièmes de l'unité sur $\mathrm{Spec}\,A$ :
\[
\mu_5^{*} = \mathrm{Spec}\,B,\qquad B = A[T]/(1 + T + T^{2} + T^{3} + T^{4}).
\]
Il est fini étale de rang $4$ sur $A$, galoisien de groupe
$(\mathbb{Z}/5\mathbb{Z})^{*} \simeq \mathbb{Z}/4\mathbb{Z}$, qui agit par
$\zeta \mapsto \zeta^{k}$ ; $2$ et $3$ en sont les générateurs, et $-1$
l'élément d'ordre $2$.\footnote{La page note ce groupe $\Gamma$ ; la lettre
est prise, en I, par le groupe de l'icosaèdre.} Soit $\beta$ l'inversion
$\zeta \mapsto \zeta^{-1}$. Dans la base $1, T, T^{2}, T^{3}$,
\[
\beta(c_0 + c_1T + c_2T^{2} + c_3T^{3}) = (c_0 - c_1) - c_1T + (c_3 - c_1)T^{2} + (c_2 - c_1)T^{3},
\]
de sorte que les invariants sont les $c_0 + c_2(T^{2} + T^{3})$ : l'anneau
$B^{\beta}$ est libre de rang $2$ sur $A$, de base $1$ et
$T^{2} + T^{3} = \zeta^{2} + \zeta^{-2}$. Ce calcul vaut sur tout anneau
$A$.\footnote{Le signe entre $c_0$ et $c_1$ est couvert d'une tache sur la
page ; la ligne suivante, que l'on suit, suppose
$c_0 - c_1(1 + T + T^{2} + T^{3})$.}

Posons $\eta = \zeta + \zeta^{-1}$ et $\eta' = \zeta^{2} + \zeta^{-2}$, son
conjugué par $\zeta \mapsto \zeta^{2}$ ; on a $\eta' = -1 - \eta$, donc
$B^{\beta} = A[\eta]$.\footnote{La page écrit ce nombre $\eta$ en tête, puis
$\gamma$ dans le calcul de la trace et de la norme ; la page 12 l'écrit
$\gamma$ et $t$, la page 21 $\gamma$ et $c$. On garde $\eta$ partout.} Sa
trace et sa norme sont
\[
\mathrm{Tr}\,\eta = \eta + \eta' = \zeta + \zeta^{2} + \zeta^{3} + \zeta^{4} = -1,
\qquad
\mathrm{N}\,\eta = \eta\,\eta' = \zeta^{3} + \zeta^{-1} + \zeta + \zeta^{-3} = -1,
\]
la seconde égalité étant celle de la page 12.

La page esquisse ensuite la situation pour $n$ quelconque : sur
$\mathbb{Z}[\frac1n]$ et pour $n \geqslant 3$, l'inversion agit librement sur
$\mu_n^{*}$, le quotient $\vartheta_n = \mu_n^{*}/(\zeta \sim \zeta^{-1})$
est fini étale de rang $\varphi(n)/2$, $\mu_n^{*} \to \vartheta_n$ est étale
de rang $2$, et $\lambda \mapsto \lambda + \lambda^{-1}$ plonge $\vartheta_n$
dans la droite affine --- ce qui revient au fait classique que l'anneau des
entiers du sous-corps réel de $\mathbb{Q}(\zeta_n)$ est
$\mathbb{Z}[\zeta_n + \zeta_n^{-1}]$.\footnote{Le rang de $\vartheta_n$ est
une lecture incertaine sur la page. La page esquisse aussi une version
« tordue » de ces objets par un torseur $T$ ($\mu_n^{T}$, $\vartheta_n^{T}$,
qu'elle note aussi $\mathcal{V}_n$), sans la développer ; on ne
l'interprète pas.}

\pagerange{12}{12}
\subsection*{L'équation (page 12)}

De la trace et de la norme :
\[
\eta^{2} + \eta - 1 = 0,\qquad \mathcal{O}(\vartheta_5) = \mathbb{Z}[\tfrac15][t]/(t^{2} + t - 1).
\]
Le discriminant du polynôme est $5$ : $\vartheta_5$ est étale sur
$\mathbb{Z}[\frac15]$, et non au-dessus de $5$.\footnote{La lettre qui nomme
cet anneau sur la page, un $\mathfrak{D}$ ou un $\vartheta$ gothique, n'est
pas sûre ; on reprend le $\vartheta_n$ de la page 14. La marge dit que la
formule $\eta = (-1 \pm \sqrt5)/2$ « garde un sens sur tout anneau où 2
\emph{et} 5 sont inversibles » : plus exactement, les deux racines sont
distinctes dès que $5$ est inversible, et la formule les exprime là où, de
plus, $2$ est inversible et $5$ a une racine carrée choisie.} Sur
$\mathbb{R}$,
\[
\eta = 2\cos\tfrac{2\pi}{5} = \frac{\sqrt5 - 1}{2},\qquad \eta' = 2\cos\tfrac{4\pi}{5} = -\frac{\sqrt5 + 1}{2} :
\]
$\eta$ est l'inverse du nombre d'or, $\eta'$ son opposé. Les deux racines
s'échangent par $t \mapsto t^{2} - 2$, trace de l'élévation au carré :
$\eta' = \eta^{2} - 2$ et $\eta = \eta'^{2} - 2$. La page note enfin
$\cos'(D, \zeta D) = \mathrm{Tr}\,\zeta^{2} = t^{2} - 2$ : pour deux droites
d'un plan, l'invariant naturel n'est pas le cosinus de leur angle, défini au
signe près, mais $\cos' = 2\cos$ de l'angle double. C'est la notation que
reprend la page 21.

\pagerange{21}{21}
\section*{IV. L'icosaèdre relativement à une forme quadratique (page 21)}

La page 21 transporte l'icosaèdre sur une base quelconque. Elle écrit
« icosaèdre projectif --- icosaèdre » : ce qui se transporte bien n'est pas
l'ensemble des douze sommets, mais celui des six droites qui joignent les
sommets opposés.

Dans l'espace euclidien de la section I, avec $Q(x) = x \cdot x$, de forme
polaire $B(x, y) = Q(x + y) - Q(x) - Q(y) = 2\,x \cdot y$ et $B' = B/2$, ces six
droites $D_1, \ldots, D_6$ font deux à deux le même angle $\alpha$ ou
$\pi - \alpha$ : ce sont six droites \emph{équiangulaires}, le maximum possible
en dimension $3$. La page le calcule à partir du nombre d'or :
\[
\cos\alpha = \frac{\eta}{2 - \eta} = \frac{2\eta + 1}{5} = \frac{1}{\sqrt5},
\qquad \cos 2\alpha = -\frac35,
\]
la forme $a + b\eta$ étant obtenue en identifiant à
$\eta^{2} + \eta - 1$.\footnote{La première équation en $a, b$ s'écrit sur
la page $2a + (2b - 1 - a)\gamma + b\gamma^{2} = 0$ ; le signe du dernier
terme doit être $-$, comme le confirme la forme normalisée qui suit, juste
sur la page, et qui donne $b = 2a$, $a = \frac15$, $b = \frac25$.} Pour $x$ et
$y$ engendrant deux droites distinctes,
\[
\cos'(D_i, D_j) = \frac{B(x,y)^{2}}{Q(x)Q(y)} - 2 = 2\cos 2\alpha = -\frac65,
\]
\[
B(x,y)^{2} = \frac45\,Q(x)Q(y),
\qquad
B'(x,y)^{2} = \frac15\,Q(x)Q(y).
\]

La page ne définit pas l'« icosaèdre local ». Ce qu'elle calcule conduit à la
définition suivante, qui est la nôtre. Soit $S$ un schéma sur lequel $2$ et
$5$ sont inversibles, $\mathcal{E}$ un $\mathcal{O}_S$-module localement
libre de rang $3$ muni d'une forme quadratique non dégénérée $Q$, et
$B' = B/2$. Un icosaèdre (projectif) relatif à $Q$ est une famille de six
sous-modules $D_1, \ldots, D_6$ localement facteurs directs de rang $1$, tels
que, pour des générateurs locaux $x_i$, les $Q(x_i)$ soient inversibles et
\[
B'(x_i, x_j)^{2} = \tfrac15\,Q(x_i)Q(x_j)\qquad (i \neq j).
\]
La page annonce alors, sans les démontrer :

\begin{itemize}
  \item[a)] il existe localement des icosaèdres ;
  \item[b)] deux icosaèdres sont localement conjugués par $\mathrm{SO}(Q)$ ;
  \item[c)] le stabilisateur d'un icosaèdre dans $\mathrm{SO}(Q)$ --- son
  groupe d'automorphismes --- est un schéma en groupes fini étale sur $S$,
  dont les fibres géométriques sont isomorphes au groupe alterné
  $\mathfrak{A}_5$.
\end{itemize}

Sur $\mathbb{R}$, c) est le théorème de la section II. On ne vérifie pas ces
énoncés en toute généralité, en particulier en caractéristique $3$, qui
divise $|\mathfrak{A}_5|$.\footnote{La page ne dit pas sur quelle base elle
se place ; l'hypothèse « $2$ et $5$ inversibles » est la nôtre, imposée par la
formule encadrée ($B' = B/2$, et le facteur $\frac15$). En tête de la page,
une colonne $\boldsymbol{\mu}_5$ --- $\boldsymbol{\mu}_5^{*}$ ---
$\mathbb{V}_5$ --- $\mathbb{Z}[\frac14]$ : la transcription lit $\frac14$,
tel qu'écrit, là où l'on attendrait $\frac15$. La page note aussi
qu'« en caractéristique $2$ », la relation en $B$ dégénère en
$B(x,y)^{2} = 0$, et passe à $B'$, qui lui-même demande $2$ inversible ; la
ligne « ceci signifie : $y \parallel x$ » qui suit reste obscure.} « Localement »
doit s'entendre pour une topologie qui adjoint $\sqrt5$ --- la topologie
étale, par exemple : sur un corps où $5$ n'est pas un carré, trois droites
$D_1, D_2, D_3$ ne peuvent être toutes rationnelles, car le produit des
$Q(x_i)Q(x_j)/5$ sur les trois paires serait à la fois un carré et $5$ fois
un carré.\footnote{Cette remarque est la nôtre ; la page ne précise pas la
topologie.}

Deux éléments de la page ne s'accordent pas avec ce calcul. En marge des
énoncés, la page écrit $c^{2} + c - 1 = 0$ pour la constante
$c = \cos'(D_i, D_j)$ ; or $c = -\frac65$ ne vérifie pas cette équation. La
relation du nombre d'or est celle de $\eta = 2\cos\frac{2\pi}{5}$, qui sert à
calculer $\cos\alpha$ mais n'est pas l'invariant d'une paire
d'axes.\footnote{On lit sans doute là une autre idée de la définition,
où la constante serait un nombre d'or, que le calcul de la même page ne
soutient pas. Le petit diagramme du haut de la page, où un schéma
$\mathrm{Ic}_Q$ est dessiné au-dessus de $\mathbb{V}_{5s}$, lui-même au-dessus
de $S$, avec une flèche $S \to \mathbb{V}_5$, appartient vraisemblablement à
cette autre idée ; $\mathbb{V}_5$ est le $\vartheta_5$ des pages 14 et 12,
que la page 14 note aussi $\mathcal{V}_n$. La page ne dit pas ce que sont
ces flèches, et l'indice $5s$ est d'une lecture douteuse ; on ne les
interprète pas.}

\pagerange{12}{15}
\section*{V. Une somme dans une catégorie d'ensembles ordonnés (pages 15, 13 et pied de la page 12)}

Une suite au crayon, qui se lit dans l'ordre : page 15, page 13, deux lignes
au pied de la page 12. Les pages 15 et 13 sont barrées chacune d'un long
trait oblique ; les deux lignes du pied de la page 12 ne le sont pas. La page
13 et ces deux lignes sont écrites tête-bêche par rapport à l'encre ; la page
15 est dans le même sens que les pages d'encre. Elle commence au milieu d'un raisonnement : la catégorie qu'elle
appelle $(\mathrm{OM})$ et les hypothèses a), b), c) de sa forme à deux
facteurs sont sur une page qui n'est pas dans ce dossier.

\pagerange{15}{15}
\subsection*{La catégorie, et les hypothèses (page 15)}

Ce que la démonstration utilise fixe la catégorie. Ses objets sont des
ensembles ordonnés ayant des bornes supérieures quelconques et des bornes
inférieures finies, les secondes distribuant sur les premières :
\[
x \wedge \bigvee_j y_j = \bigvee_j\, (x \wedge y_j) ;
\]
ses morphismes sont les applications qui commutent aux bornes supérieures
quelconques et aux bornes inférieures finies. Ce sont les \emph{frames} :
l'ensemble ordonné des ouverts d'un espace topologique en est le modèle, et
leur catégorie opposée est celle des \emph{locales}.\footnote{Le sigle
$(\mathrm{OM})$ n'est pas expliqué dans ce dossier, et la lecture des deux
lettres n'est pas sûre. L'identification avec les frames est la nôtre :
elle résulte de ce que les démonstrations des pages 15 et 13 utilisent,
y compris la distributivité dans l'objet d'arrivée, que la page emploie sans
la dire. Les lettres $U$, $V$, $W$ pour les éléments, et l'écriture
« $\gamma(U \times V)$ », suggèrent le modèle des ouverts d'un produit. Les
frames, sous d'autres noms, remontent à Ehresmann (1957) et Bénabou (1958) ;
le mot \emph{frame} est de Dowker, celui de \emph{locale} d'Isbell (1972),
et les sommes de frames ont été étudiées dans les années 1970, notamment
par Dowker et Strauss. Rien sur ces pages ne cite ces travaux.}

Soient $I$, $J$, $K$ trois frames et $m : I \times J \to K$ une application ;
on pose $a(U) = m(U, 1)$ et $b(V) = m(1, V)$.\footnote{La page écrit
$\gamma(U \times V)$, $\alpha$ et $\beta$ ; on écrit $m(U,V)$, $a$ et $b$, les
lettres $\alpha$, $\beta$, $\gamma$ étant prises ailleurs dans le dossier. De
même, $f_1$, $f_2$ et $f$ ci-dessous sont les $\bar\alpha$, $\bar\beta$ et
$\bar\gamma$ de la page.} Les hypothèses dont la démonstration se sert sont
les suivantes :

\begin{itemize}
  \item[a)] $a$ et $b$ commutent aux bornes supérieures quelconques ;
  \item[b)] $m$ commute aux bornes inférieures finies :
  $m(1,1) = 1$ et $m(U \wedge U', V \wedge V') = m(U,V) \wedge m(U',V')$ ; en
  particulier $m(U,V) = a(U) \wedge b(V)$, et $a$, $b$ sont des morphismes ;
  \item[c)] tout élément de $K$ est borne supérieure d'éléments $m(U,V)$ ;
  \item[d)] si $m(U,V) \leqslant a(U')$, alors $U \leqslant U'$ ou $V = 0$ ; et
  symétriquement pour $b$ ;
  \item[e)] si $m(U,V) \leqslant \bigvee_i W_i$, il existe des décompositions
  $U = \bigvee_j U_j$ et $V = \bigvee_k V_k$ telles que chaque $m(U_j, V_k)$
  soit majoré par l'un des $W_i$.
\end{itemize}

Les hypothèses a), b), c) sont reconstituées d'après leur usage. Les
hypothèses d) et e) sont écrites au pied de la page 12, mais sous la forme
infinie de la page 13 ; la page 15 invoque d) sous la forme à deux facteurs
donnée ici, et la page 13 utilise e) sous cette forme sans la
nommer.\footnote{La forme d) du pied de la page 12 est :
« $\gamma((V_\lambda)) \leq \alpha_\lambda(V'_\lambda) \Rightarrow V_\lambda \leq V'_\lambda$
ou $\exists\,\mu \neq \lambda$ avec $V_\mu = 0$ ». Elle se spécialise
exactement en la forme à deux facteurs, et c'est celle que la page 15
utilise.}

\textbf{Proposition.} Sous ces hypothèses, $(a, b)$ fait de $K$ la somme de
$I$ et $J$ dans la catégorie des frames.

\emph{Unicité.} Si $\varphi : K \to L$ est un morphisme, il est connu sur les
$a(U)$ et les $b(V)$, donc sur leurs bornes inférieures
$m(U,V) = a(U) \wedge b(V)$ par b), donc sur leurs bornes supérieures, donc
sur tout $K$ par c).

\emph{Existence.} Soient $f_1 : I \to L$ et $f_2 : J \to L$ deux morphismes
vers un frame $L$. On pose
\[
f(W) = \bigvee_{m(U,V) \leqslant W} f_1(U) \wedge f_2(V).
\]
1°) $f \circ a = f_1$. Les couples $(U, V)$ avec $m(U,V) \leqslant a(U')$
sont, par d), ceux pour lesquels $V = 0$, dont le terme est nul, et ceux pour
lesquels $U \leqslant U'$ ; la borne supérieure vaut donc
$f_1(U') \wedge f_2(1) = f_1(U')$. De même $f \circ b = f_2$.\footnote{La page
écrit « ou bien $V = 0_I$ (l'élément correspondant du Sup est alors $0_K$) » :
il faut $0_J$, et le terme nul est $0_L$. Plus bas, elle écrit
$\bar\gamma(1_K) = \ldots = 1_K$, pour $1_L$.}

2°) $f$ est un morphisme. Elle est croissante ; $f(1) = f_1(1) \wedge f_2(1) = 1$.
Pour les bornes inférieures finies, b) donne
$m(U \wedge U', V \wedge V') \leqslant W \wedge W'$ dès que
$m(U,V) \leqslant W$ et $m(U',V') \leqslant W'$, d'où, par la distributivité
dans $L$,
\[
f(W \wedge W') = \bigvee_{\substack{m(U,V) \leqslant W \\ m(U',V') \leqslant W'}} \bigl(f_1(U) \wedge f_2(V)\bigr) \wedge \bigl(f_1(U') \wedge f_2(V')\bigr) = f(W) \wedge f(W').
\]
Pour les bornes supérieures : $f(\bigvee W_i) \geqslant \bigvee f(W_i)$ est
clair. Inversement, si $m(U,V) \leqslant \bigvee W_i$, on décompose par e) ;
chaque $f_1(U_j) \wedge f_2(V_k)$ est majoré par un $f(W_i)$, et
$f_1(U) \wedge f_2(V) = \bigvee_{j,k} f_1(U_j) \wedge f_2(V_k)$ par
distributivité.

L'hypothèse e) demande une subdivision « en damier », unique pour tout le
recouvrement. Dans la somme de deux frames, la relation de recouvrement
s'obtient par subdivisions successives, qui ne se ramènent pas a priori à un
damier : e) est donc une vraie hypothèse, et la proposition un critère
suffisant.\footnote{Cette remarque est la nôtre ; la page ne discute pas des
cas où les hypothèses sont satisfaites. Pour les frames d'ouverts, on sait
aujourd'hui que le frame des ouverts de $X \times Y$ est la somme de ceux de
$X$ et de $Y$ lorsque $X$ est localement compact, mais pas en général.}

\pagerange{13}{13}
\subsection*{Sommes infinies (page 13 et pied de la page 12)}

La page généralise à une famille $(I_\lambda)_{\lambda \in \Lambda}$. Le
produit restreint $\prod' I_\lambda$ est formé des familles $(V_\lambda)$ avec
$V_\lambda = 1$ pour presque tout $\lambda$, et
$\alpha^{0}_\mu : I_\mu \to \prod' I_\lambda$ insère un élément à la place
$\mu$ et $1$ ailleurs. Pour $m : \prod' I_\lambda \to K$ et
$a_\mu = m \circ \alpha^{0}_\mu$, les hypothèses sont : a) les $a_\mu$
commutent aux bornes supérieures quelconques ; b) $m$ commute aux bornes
inférieures finies, de sorte que $m((V_\lambda)) = \bigwedge_\lambda a_\lambda(V_\lambda)$,
borne inférieure en fait finie ; c) tout élément de $K$ est borne supérieure
d'éléments de l'image de $m$ ; d) et e), au pied de la page 12, comme
ci-dessus.\footnote{La page écrit, au c) : « tout élément de $K$ est le sup
d'éléments de $\prod'_\lambda I_\lambda$ » ; il faut « de l'image de
$\prod' I_\lambda$ par $m$ ». Elle s'achève sur « TSVP » ; la suite est aux
deux lignes du pied de la page 12, qui énoncent d) et le début de e), et
s'arrêtent.} La conclusion attendue --- $K$ est la somme des $I_\lambda$ ---
n'est pas écrite ; la démonstration à deux facteurs s'y transpose mot pour
mot.

\pagerange{22}{22}
\section*{VI. Le produit d'un topos par un topos de préfaisceaux (page 22)}

Une page barrée d'une croix, entièrement lisible. Soient $E$ un topos de
Grothendieck et $C$ une petite catégorie. La catégorie
\[
\pi(C, E) = \underline{\mathrm{Hom}}(C^{\circ}, E)
\]
des préfaisceaux sur $C$ à valeurs dans $E$ est un topos. Un foncteur
$f : C \to C'$ et un morphisme de topos $g : E \to E'$ définissent un
morphisme de topos $\pi(f, g) : \pi(C, E) \to \pi(C', E')$, d'image inverse
et d'image directe
\[
\pi(f,g)^{*}(u') = g^{*} \circ u' \circ f^{\circ},
\qquad
\pi(f,g)_{*}(u) = g_{*} \circ f_{*}(u),
\]
où $f_{*}$ est l'extension de Kan à droite le long de $f^{\circ}$, calculée
dans $E$ ; comme $g_{*}$ commute aux limites, on peut aussi la calculer après
$g_{*}$, dans $E'$. Ainsi $\pi$ est un $2$-foncteur
$(\mathrm{Cat}) \times (\mathrm{Top}) \to (\mathrm{Top})$.\footnote{La page
écrit $\pi(f,g)_{*} = (u' \mapsto g_{*} f_{*}^{(E)})$ et note en marge
« vérifier l'expression donnée en $\pi(f,g)_{*}$ » : l'expression est juste,
appliquée à $u \in \pi(C, E)$ plutôt qu'à $u'$.}

On définit un morphisme de topos
\[
(1)\qquad \gamma(C, E) : \pi(C, E) \longrightarrow \hat{C} \times E
\]
vers le produit dans la $2$-catégorie des topos, par ses deux composantes :
$\gamma_1 = \pi(\mathrm{id}_C, p_E)$, où $p_E : E \to P$ est le morphisme vers
le topos ponctuel et $\hat{C} = \pi(C, P)$, d'image inverse
$F \mapsto (X \mapsto F(X)_E)$, le préfaisceau d'objets constants ; et
$\gamma_2 = \pi(p_C, \mathrm{id}_E)$, où $p_C : C \to C_0$ est le foncteur vers
la catégorie ponctuelle et $E = \pi(C_0, E)$, d'image inverse
$X \mapsto (i \mapsto X)$, le préfaisceau constant.\footnote{La page qualifie
$p_C$ d'« unique morphisme de topos », d'une lecture incertaine : c'est un
foncteur.}

\textbf{Question} (de la page). $\gamma(C, E)$ est-il toujours une
équivalence ? La page note trois cas. a) Les deux membres transforment
sommes en sommes, en chacun des arguments ; on peut donc supposer $C$
connexe. b) Si $C$ est un groupoïde connexe, équivalent à un groupe $G$,
alors $\hat{C}$ est le topos classifiant $B_G$ et $\pi(C, E)$ la catégorie
des objets de $E$ munis d'une action de $G$. c) Si $E = \hat{D}$, alors
$\pi(C, E) \simeq (C \times D)^{\wedge}$, et $\gamma$ est le morphisme évident
$(C \times D)^{\wedge} \to \hat{C} \times \hat{D}$.

La réponse est oui : que la catégorie des préfaisceaux sur $C$ à valeurs dans
un topos de Grothendieck $E$ soit le produit de $\hat{C}$ et de $E$ --- ce
qu'on énonce aussi en disant que le changement de base de $\hat{C}$ le long
de $E \to P$ est $\underline{\mathrm{Hom}}(C^{\circ}, E)$ --- est aujourd'hui
un résultat classique de la théorie des topos (voir par exemple
P. T. Johnstone, \emph{Sketches of an Elephant}, 2002).\footnote{Le dossier
161-4 du même fonds travaille la même comparaison, avec les mêmes notations
$\pi(C, E)$ et $\gamma(C, E)$, aux pages 15 et 17 de sa transcription.
Page 17, il établit le cas c) quand $C$ et $D$ ont des limites projectives
finies, par la propriété universelle d'un topos de préfaisceaux
($\underline{\mathrm{Homtop}}(E, \hat{A}) \simeq \underline{\mathrm{Hom}}^{*}_{\mathrm{gex}}(A, E)^{\circ}$
pour $A$ à limites finies) ; puis, pour $E = \widetilde{D}$ quelconque, il
plonge $C$ dans une catégorie à limites finies $C'$ et $E$ dans $\hat{D}$,
montre que $\gamma(C, E)$ est un \emph{plongement}, et s'interrompt sur la
phrase « ce qui prouvera que $\gamma$ est une équivalence ». La page 15 de
161-4 est un fragment de la même démonstration. Les deux dossiers ne sont pas
fusionnés ici ; on ne sait pas lequel précède l'autre.}

\pagerange{24}{30}
\section*{VII. Graphes plongés dans une surface orientée (pages 24 à 30)}

\pagerange{24}{24}
\subsection*{Les ordres circulaires aux sommets (page 24)}

Soit $X$ un espace compact triangulable de dimension au plus $1$ --- un
graphe fini --- et $S$ une surface orientée. On étudie les classes
d'isotopie de plongements modérés $\varphi : X \hookrightarrow S$ ; leur
ensemble est noté $\mathrm{Pl}_{\mathrm{isot}}(X, S)$.

Les sommets topologiques de $X$ sont les points où $X$ n'est pas localement
un intervalle : points isolés, bouts (une seule branche), points de
ramification. On note $X'_0$ l'ensemble de ces derniers, où partent
$\nu_x \geqslant 3$ branches. Un plongement induit en chaque $x \in X'_0$ un
ordre circulaire $\omega_x(\varphi)$ sur les branches issues de $x$ --- celui
dans lequel on les rencontre en tournant autour de $\varphi(x)$ dans le sens
de l'orientation ; l'ensemble $\Omega_x(X)$ de ces ordres a
$(\nu_x - 1)!$ éléments.\footnote{La page parle des « arêtes issues de
$x$ » ; si une arête est une boucle en $x$, elle compte pour deux branches,
et c'est sur les branches qu'il faut prendre l'ordre. Aux bouts, la donnée
est vide ; aux points de valence $2$, qui ne sont pas des sommets
topologiques, elle est unique.} On obtient une application
\[
\omega : \mathrm{Pl}_{\mathrm{isot}}(X, S) \longrightarrow \Omega(X) = \prod_{x \in X'_0} \Omega_x(X),
\]
et l'on se propose d'étudier ses fibres $\mathrm{Pl}^{\alpha}_{\mathrm{isot}}(X, S)$,
ou plus généralement, pour $X$ et $\alpha \in \Omega(X)$ fixés, tous les
plongements $\varphi$ de $X$ dans des surfaces orientées tels que
$\omega(\varphi) = \alpha$.

Dans le langage d'aujourd'hui, un couple $(X, \alpha)$ est un \emph{graphe
enrubanné} (\emph{ribbon graph}, \emph{fat graph}), et $\alpha$ un
\emph{système de rotation}.\footnote{La notion de système de rotation remonte
à Heffter (1891) et a été formalisée par Edmonds (1960) ; elle est au cœur
de la démonstration du théorème de coloriage des cartes (Ringel et Youngs,
1968). Les termes \emph{ribbon graph} et \emph{fat graph} datent
des années 1980. Rien sur la page ne cite ces travaux. Les graphes
enrubannés sont aussi des objets que l'\emph{Esquisse d'un programme} (1984)
étudie, sous le nom de « cartes » ; ces pages ne la citent pas. Signalons enfin, sans pouvoir dire
s'il y a un lien, que le titre d'un article d'Y. Ladegaillerie,
« Classes d'isotopie de plongements de 1-complexes dans les surfaces »
(\emph{Topology}, 1984), issu d'un travail fait à Montpellier dans les années
1970, décrit exactement le sujet de ces pages ; rien sur elles ne le
mentionne.}

\pagerange{24}{25}
\subsection*{Le voisinage tubulaire, et ses bords (pages 24 et 25)}

\textbf{Théorème 1.} Pour $\alpha$ fixé, le voisinage tubulaire fermé d'un
plongement $\varphi$ avec $\omega(\varphi) = \alpha$ ne dépend, à un
homéomorphisme près induisant l'identité sur $X$ --- unique à isotopie
près ---, que de $\alpha$, et non de $S$ ; on le note $\mathcal{U}_\alpha$.
La restriction à $X$ induit une bijection
\[
\mathrm{Pl}^{+}_{\mathrm{isot}}(\mathcal{U}_\alpha, S) \xrightarrow{\ \sim\ } \mathrm{Pl}^{\alpha}_{\mathrm{isot}}(X, S),
\]
où l'exposant $+$ désigne les plongements qui préservent
l'orientation.\footnote{La page écrit la flèche dans l'autre sens,
$\mathrm{Pl}^{\alpha}_{\mathrm{isot}}(X,S) \to \mathrm{Pl}_{\mathrm{isot}}(\mathcal{U}_\alpha, S)$
(par passage au voisinage tubulaire), et « celle-ci est une bijection », ce
dernier mot d'une lecture incertaine. La précision « préservant
l'orientation » est la nôtre : un plongement de $\mathcal{U}_\alpha$ qui
renverse l'orientation induirait sur $X$ l'ordre opposé à $\alpha$. La page
ne démontre pas le théorème.}

$\mathcal{U}_\alpha$ est une surface compacte orientée à bord, qui se
rétracte par déformation sur $X$ : c'est la surface-ruban du graphe
enrubanné. Son complémentaire
$\mathcal{V}_\alpha = \mathcal{U}_\alpha \smallsetminus X$ a pour composantes
connexes des anneaux semi-ouverts $\mathcal{V}^{i}_\alpha \simeq S^{1} \times [1, 2[$,
dont l'extrémité fermée $S^{1} \times \{1\}$ est une composante du bord de
$\mathcal{U}_\alpha$, et dont l'extrémité ouverte est collée à $X$ : les
voisinages de $X$ dans $\mathcal{U}_\alpha$ induisent sur
$\mathcal{V}^{i}_\alpha$ le filtre des $S^{1} \times {]2 - \varepsilon, 2[}$.\footnote{Telle
que la transcription la lit, la page donne à la fois
$\mathrm{bord}(\mathcal{U}_\alpha) \cap \mathcal{V}^{i}_\alpha \simeq S^{1} \times \{1\}$
et, pour trace des voisinages de $X$, « le filtre des voisinages de
$S^{1} \times \{1\}$ dans $S^{1} \times {]0,2[}$ » : les deux extrémités
seraient la même. Les chiffres sont récrits sur d'autres et le passage est
lourdement raturé. On donne l'énoncé cohérent.}

On note $I(\alpha) = \pi_0(\mathcal{V}_\alpha)$ l'ensemble de ces anneaux,
c'est-à-dire des composantes du bord de $\mathcal{U}_\alpha$. La page demande
de l'« expliciter combinatoirement » ; voici comment.\footnote{Cette
description et la formule qui suit sont les nôtres.} Soit $\mathsf{H}$
l'ensemble des demi-arêtes de $X$ (pour une triangulation quelconque), $\iota$
l'involution qui échange les deux moitiés d'une arête, et $\sigma$ la
permutation qui envoie une demi-arête sur la suivante autour de son sommet
pour l'ordre $\alpha$. Les composantes du bord de $\mathcal{U}_\alpha$
correspondent aux orbites de $\sigma \circ \iota$, auxquelles s'ajoute un
cercle par point isolé de $X$. Si $g(\alpha)$ désigne la somme des genres des
surfaces fermées obtenues en bouchant chaque bord de chaque composante de
$\mathcal{U}_\alpha$ par un disque, la formule d'Euler donne
\[
|I(\alpha)| = b_0(X) + b_1(X) - 2\,g(\alpha).
\]
La marge de la page 25 conjecture « le nombre des $\mathcal{V}^{i}_\alpha$ est
$b_1(X) + 1$ si $X$ est connexe non vide », et en général
$b_1(X) + b_0(X)$ : c'est exact si et seulement si $g(\alpha) = 0$.

\pagerange{25}{26}
\subsection*{Recoller le complémentaire : les surfaces modèles (pages 25 et 26)}

Pour un plongement de $\mathcal{U}_\alpha$ dans $S$, la surface $S$ s'obtient
en recollant à $\mathcal{U}_\alpha$ la surface orientée à bord
$\mathcal{U}' = \overline{S \smallsetminus \varphi(\mathcal{U}_\alpha)}$,
compacte si et seulement si $S$ l'est, par un isomorphisme de
$\partial\,\mathcal{U}'$ sur $-\partial\,\mathcal{U}_\alpha \simeq I(\alpha) \times S^{1}$.
Si toute composante connexe de $S$ rencontre $X$, toute composante connexe de
$\mathcal{U}'$ a un bord non vide. Trouver tous les plongements de $X$ dans
une surface orientée, à isomorphisme près induisant l'identité sur $X$,
revient donc à trouver toutes les façons d'écrire $I(\alpha) \times S^{1}$
comme bord orienté d'une surface.

Une telle surface est décrite par a) une partition $\pi$ de $I = I(\alpha)$
en parties $(I_\lambda)_{\lambda \in \Lambda}$ ; b) pour chaque $\lambda$,
une surface orientée connexe $\mathcal{U}'_\lambda$ et un isomorphisme
$\partial\,\mathcal{U}'_\lambda \simeq -I_\lambda \times S^{1}$. Si l'on se
borne aux surfaces « algébriques », compactes moins un ensemble fini, une
telle $\mathcal{U}'_\lambda$ est déterminée par son genre $g_\lambda$ --- celui
de la surface fermée obtenue en rajoutant les bouts et en contractant chaque
composante du bord en un point --- et par son nombre de bouts
$\nu_\lambda$.\footnote{La page 26 écrit le nombre de bouts d'une lettre que
la transcription lit $\beta$, et la page 27 d'un $\nu$ net ; c'est la même
lettre. On écrit $\nu$.} Explicitement,
\[
\mathcal{U}'_\lambda = S_{g_\lambda} \smallsetminus \Bigl(\bigcup_{j \in I_\lambda} \mathring{D}_j \cup P_\lambda\Bigr),
\]
où $S_g$ est la surface fermée orientée de genre $g$, les $D_j$ des disques
fermés disjoints, et $P_\lambda$ un ensemble de $\nu_\lambda$ points hors des
disques. Un homéomorphisme orienté de $\mathcal{U}'_\lambda$ peut permuter
les composantes du bord de façon prescrite ; il est donc inutile de préciser
laquelle est recollée à quel $j$.

Soit $\rho = (\alpha, \pi, (g_\lambda), (\nu_\lambda))$ et
$S_\rho = \mathcal{U}_\alpha \cup \coprod_\lambda \mathcal{U}'_\lambda$ la
surface ainsi construite, munie d'un plongement canonique
$\varphi_\rho : X \hookrightarrow S_\rho$. Elle a
$\nu = \sum_\lambda \nu_\lambda$ bouts. Son genre, lorsqu'elle est connexe, est
\[
g(S_\rho) = \sum_{\lambda \in \Lambda} g_\lambda + g(\alpha) + b_1(\Delta_\rho),
\]
où $\Delta_\rho$ est le graphe biparti dont les sommets sont les composantes
connexes de $X$ et les parties $I_\lambda$, et dont les arêtes sont les
éléments de $I$, chacun joignant la composante de $X$ qu'il borde à la partie
qui le contient.\footnote{La marge de la page 26 propose
« $\sum_\lambda g_\lambda + \gamma(\alpha)$, $\gamma(\alpha)$ un entier (le
genre de $\mathcal{U}_\alpha$) --- \emph{vérifier} ». La formule est juste
exactement lorsque $\Delta_\rho$ est un arbre ; en général il faut ajouter
$b_1(\Delta_\rho)$, qui compte les anses créées par le recollement
lui-même. Exemple : $X$ un cercle, $\mathcal{U}_\alpha$ un anneau, les deux
bords recollés à un même anneau $\mathcal{U}'$ ; on obtient un tore, alors
que $g_\lambda = g(\alpha) = 0$. La formule donnée ici, qui se déduit de la
caractéristique d'Euler, est la nôtre.}

\textbf{Théorème 2.} Tout plongement de $X$ dans une surface orientée
algébrique est isomorphe à un $\varphi_\rho$, pour un $\rho$ unique à
renumérotation des $\lambda$ près.\footnote{La phrase de la page s'arrête au
bas de la page 26 sur « uniquement déterminés … ». La précision « à
renumérotation près » est la nôtre.}

\pagerange{27}{28}
\subsection*{Le groupe modulaire (pages 27 et 28)}

Pour $S$ fixée, on a donc une application de $\mathrm{Pl}_{\mathrm{isot}}(X, S)$
vers l'ensemble des $\rho$, et l'on est ramené à étudier, pour $\rho$ fixé, la
fibre $\mathrm{Pl}_{\mathrm{isot}}(X, S; \rho)$, non vide si et seulement si
$S \simeq S_\rho$. La restriction à $X$ définit une application
\[
(*)\qquad r : \mathrm{Isot}^{+}(S_\rho, S) \longrightarrow \mathrm{Pl}_{\mathrm{isot}}(X, S; \rho),
\]
surjective d'après le Théorème 2, où $\mathrm{Isot}^{+}$ désigne les classes
d'isotopie d'homéomorphismes préservant l'orientation.

Soit $\mathrm{Mod}(S)$ le groupe des classes d'isotopie d'homéomorphismes de
$S$ préservant l'orientation --- le \emph{groupe modulaire} de $S$, qu'on
appelle aussi \emph{mapping class group}, et que la page appelle « groupe de
Teichmüller ».\footnote{La page note ce groupe $G = \mathrm{Isot}_o(S,S)$,
puis $G = \mathrm{Isot}(S,S)$ ; elle note aussi $G_\rho$ à la fois
$\pi_1(S_\rho)$ et $\mathrm{Isot}_o(S_\rho, S_\rho)$. On écrit
$\mathrm{Mod}(S)$ et $\Pi = \pi_1(S_\rho)$. « Groupe modulaire de
Teichmüller » était le nom courant dans les années 1970.} Il agit sur les deux
membres de $(*)$, et $r$ est équivariante ; le premier membre est un torseur
sous $\mathrm{Mod}(S)$. Choisissant $f : S_\rho \xrightarrow{\sim} S$, et
notant $H \subset \mathrm{Mod}(S)$ le stabilisateur de la classe de
$f|_X$, l'application $r$ s'identifie à
$\mathrm{Mod}(S) \to \mathrm{Mod}(S)/H$. Il reste deux questions : déterminer
$\mathrm{Mod}(S)$ ; déterminer $H$, le sous-groupe des $h$ tels que
$h \circ f|_X$ soit isotope à $f|_X$. On peut supposer $S = S_\rho$ et
$f = \varphi_\rho$.

La première est classique, et c'est le théorème de Dehn, Nielsen et Baer :
pour $S$ fermée de genre au moins $1$, l'homomorphisme canonique
\[
\mathrm{Mod}(S) \longrightarrow \mathrm{Out}(\pi_1(S, s))
\]
est injectif, d'image le sous-groupe d'indice $2$ des automorphismes
extérieurs qui préservent l'orientation ; il est injectif aussi, plus
généralement, pour les surfaces algébriques.\footnote{La page 28 écrit que
cet homomorphisme « est un isomorphisme », et la page 27 que
$\mathrm{Isot}(S_\rho, S)$ est un torseur sous
$\mathrm{Autext}(\pi_1(S_\rho))$, « du moins si $S_\rho$ est compacte ». Ce
n'est vrai que si l'on admet les homéomorphismes qui renversent
l'orientation ; or ceux-ci ne préservent pas la fibre au-dessus de $\rho$,
puisqu'ils renversent $\alpha$. Pour les classes qui préservent
l'orientation, qui sont celles qu'il faut ici, l'image est d'indice $2$. Le
genre $0$ est à part : $\mathrm{Mod}(S^{2})$ est trivial, et
$\mathrm{Out}(1)$ aussi.}

\pagerange{28}{29}
\subsection*{Le cas normal, la sphère, et une question (pages 28 et 29)}

Supposons $X$ connexe et $\pi_1(X) \to \pi_1(S)$ surjectif : la page appelle
un tel plongement \emph{normal}. Cela signifie que la partition $\pi$ est
discrète, que les $g_\lambda$ sont nuls et que chaque $\nu_\lambda$ vaut $0$
ou $1$ : $S$ se déduit de $\mathcal{U}_\alpha$ en bouchant chaque bord par un
disque, éventuellement épointé.\footnote{La page écrit, « sauf erreur »,
que les $\nu_\lambda$ sont nuls. Un disque épointé recollé le long de son
bord ne change pas le groupe fondamental, puisque la surface se rétracte
alors sur $\mathcal{U}_\alpha$ ; la condition juste est
$\nu_\lambda \leqslant 1$. Elle se réduit à $\nu_\lambda = 0$ quand $S$ est
compacte.} Le genre de $S$ est alors $g(\alpha)$.

Si $h\varphi$ est isotope à $\varphi$, il lui est homotope, donc
$\pi_1(h) \circ c$ et $c$ coïncident à automorphisme intérieur près, où
$c = \pi_1(\varphi) : \pi_1(X) \to \pi_1(S)$ ; $c$ étant surjectif,
$\pi_1(h)$ est intérieur. Par l'injectivité de Dehn--Nielsen--Baer, $h$ est
isotope à l'identité : $H = \{1\}$, et les plongements normaux de type $\rho$
forment un torseur sous $\mathrm{Mod}(S)$.\footnote{La page s'arrête à
« $\pi_1(h)$ est un automorphisme intérieur, $c$ étant surjectif, en gros » ;
la conclusion $H = \{1\}$ est la nôtre.}

En particulier, pour $S = S^{2}$, dont tout homéomorphisme orienté est
isotope à l'identité, et $X$ connexe : les classes d'isotopie de plongements
de $X$ dans la sphère correspondent bijectivement aux $\alpha \in \Omega(X)$
tels que $g(\alpha) = 0$.\footnote{L'exposant de $S^{2}$, aux deux
occurrences, est tracé comme un $1$ ; c'est la sphère que l'argument
demande, et la transcription le lit ainsi en le marquant douteux. La
phrase commencée au bas de la page 28 se ferme en tête de la page 29.} La
page 29 ajoute qu'« il faut donner une certaine condition simple sur
$\alpha$ pour qu'il en soit ainsi ». La formule de la page 25 la donne :
$g(\alpha) = 0$ si et seulement si $\sigma \circ \iota$ a exactement
$b_1(X) + 1$ orbites.\footnote{Cette réponse est la nôtre. C'est la forme
combinatoire classique du critère de planarité par les systèmes de
rotation.} C'est le cas de l'icosaèdre muni d'un système cohérent au sens
de la page 10 : $b_1 = 30 - 12 + 1 = 19$, et vingt faces.

\textbf{Question} (de la page 29). Soit $h \in \mathrm{Mod}(S)$ tel que, pour
toute composante connexe $X_i$ de $X$, de plongement induit
$\varphi_i : X_i \to S$, on ait
$\pi_1(h\varphi_i) = \pi_1(\varphi_i)$ dans
$\mathrm{Homext}(\pi_1(X_i), \pi_1(S))$ --- condition évidemment nécessaire
pour que $h \in H$. A-t-on $h \in H$ ? Autrement dit, pour ces plongements,
l'isotopie se lit-elle sur l'homotopie ? La page ne répond pas, et nous ne
tranchons pas ici.\footnote{Pour une seule courbe fermée simple, la réponse
est le théorème classique de Baer et Epstein : deux courbes simples
homotopes sont isotopes. Pour un $1$-complexe quelconque, la question
appelle des précautions (courbes parallèles, composantes contractiles) ; le
lecteur qui voudrait savoir si elle a reçu une réponse commencera par
l'article de Ladegaillerie cité plus haut.}

\pagerange{30}{30}
\subsection*{Une suite exacte (page 30)}

La dernière page de la suite écrit, pour le groupe topologique
$\mathcal{G}$ des homéomorphismes de $S$ et le sous-groupe
$\mathcal{H} = \mathrm{Aut}(S/X)$ de ceux qui induisent l'identité sur $X$,
l'identification de $\mathcal{G}/\mathcal{H}$ à un espace $\mathfrak{X}$ de
plongements --- l'orbite de $\varphi$ --- et la suite d'homotopie
\[
\pi_1(\mathcal{H}) \to \pi_1(\mathcal{G}) \to \pi_1(\mathfrak{X}) \to \pi_0(\mathcal{H}) \to \pi_0(\mathcal{G}) \to \pi_0(\mathfrak{X}).
\]
La page la marque « ?? ». Elle vaut dès que $\mathcal{G} \to \mathfrak{X}$ est
une fibration localement triviale, c'est-à-dire sous un théorème d'extension
des isotopies pour les plongements modérés de $1$-complexes, que la page ne
démontre pas.\footnote{Cette lecture est la nôtre. Les lettres de la page
sont une ronde, une cursive et un $X$ à double jambage ; que $\mathcal{G}$
soit le groupe des homéomorphismes préservant l'orientation n'est pas
précisé.} Son intérêt est de donner la réponse de principe à la seconde
question de la page 27 : l'exactitude en $\pi_0(\mathcal{G}) = \mathrm{Mod}(S)$
dit que $H$ est l'image du groupe des classes d'homéomorphismes de $S$ qui
fixent $X$ point par point.

\pagerange{32}{33}
\section*{VIII. Graphes à groupe transitif sur les sommets et les arêtes (pages 32 et 33)}

Un \emph{graphe} est ici la donnée d'un ensemble $S$ de sommets, d'un
ensemble $A$ d'arêtes orientées muni d'une involution sans point fixe
$\sigma$ --- le retournement ---, d'une partie $A' \subset A$ stable par
$\sigma$, formée des arêtes qui ont des extrémités, et d'une application
« origine » $\varphi_1 : A' \to S$, d'où
$\varphi = (\varphi_1, \varphi_1 \circ \sigma) : A' \to S \times S$.\footnote{La
page écrit
$(S, A, A' \xrightarrow{\varphi} S \times S, \sigma)$ sans définir $A'$ ;
qu'il s'agisse des arêtes non isolées est notre lecture, imposée par le cas
a), où $A' = \emptyset$. La page esquisse aussi, avant de la biffer, une
autre présentation par une relation d'incidence, et une présentation
équivalente $(S, R \subset A \times A, L \to A, \sigma)$, sans la développer.}
Une arête est \emph{circulaire} --- une boucle --- si ses deux extrémités
coïncident.

Soit $G$ un groupe agissant sur un graphe $\Gamma$, c'est-à-dire un
sous-groupe de $\mathrm{Perm}(S) \times \mathrm{Perm}(A)$ respectant la
structure, et supposons $G$ transitif sur $S$ et sur $A$.\footnote{La page
écrit « $G \simeq \mathrm{Perm}(S) \times \mathrm{Perm}(A)$ » ; c'est une
inclusion.} Comme $A'$ est stable par $G$, ou bien $A' = \emptyset$, ou bien
$A' = A$ et $\varphi_1$ est surjective ; dans ce second cas, par
transitivité, ou bien toutes les arêtes sont des boucles, ou bien aucune.

\emph{a) Les arêtes n'ont pas d'extrémités.} Alors
$\mathrm{Aut}\,\Gamma = \mathrm{Perm}(S) \times \mathrm{Aut}(A, \sigma)$, et, si
$A_1 = A/\sigma$,
$\mathrm{Aut}(A, \sigma) \simeq (\mathbb{Z}/2\mathbb{Z})^{A_1} \rtimes \mathrm{Perm}(A_1)$,
le groupe hyperoctaédral.

Dans les deux autres cas, on choisit $a \in A$ et $s = \varphi_1(a)$, et l'on
note $H = G_a \subset H' = G_s$ leurs stabilisateurs, de sorte que
$A \simeq G/H$ et $S \simeq G/H'$. Par les deux lemmes des pages 16 et 19,
l'involution $\sigma$, qui commute à $G$, est la multiplication à droite par
un élément $\dot\sigma \in N_G(H)/H$ ; elle est sans point fixe si et
seulement si $\dot\sigma \neq 1$, et involutive si et seulement si
$\dot\sigma^{2} = 1$. Le groupe $G$ agit fidèlement sur $\Gamma$ si et
seulement si $\bigcap_{g \in G} g H g^{-1} = \{1\}$.

\emph{b) Toutes les arêtes sont des boucles.} La compatibilité
$\varphi_1 \circ \sigma = \varphi_1$ signifie $\dot\sigma \in H'/H$, d'où les
données
\[
G \supset H' \supset H,\qquad \dot\sigma \in \bigl(H' \cap N_G(H)\bigr)/H,\qquad \dot\sigma^{2} = 1,\ \dot\sigma \neq 1 .
\]
Ce cas est biffé sur la page, qui reste lisible.

\emph{c) Aucune arête n'est une boucle} (page 33). Les données sont
\[
G \supset H' \supset H,\qquad \dot\sigma \in N_G(H)/H,\qquad \dot\sigma^{2} = 1,\qquad \dot\sigma \notin H'/H,
\]
la dernière condition disant que l'extrémité $\varphi_1(\sigma a)$ diffère de
l'origine.\footnote{Sous $H'$ et $H$, la page 33 inscrit $G_a$ et $G_s$,
l'inverse de la page 32 ; c'est un lapsus, car l'origine
$\varphi_1 : A \to S$ étant équivariante, $G_a \subset G_{\varphi_1(a)} = G_s$,
et la suite de la page suppose bien $H = G_a$. La condition
$\dot\sigma \notin H'/H$ n'est pas écrite ; elle est la nôtre, et c'est elle
qui sépare le cas c) du cas b).} Le graphe est \emph{simplicial} --- sans
arêtes multiples, donné par un ensemble de parties à deux éléments de $S$ ---
si et seulement si $\varphi : A \to S \times S$, c'est-à-dire
$gH \mapsto (gH', g\dot\sigma H')$, est injective, ce qui équivaut à
\[
H' \cap \dot\sigma H' \dot\sigma^{-1} \subset H,
\]
condition indépendante du représentant de $\dot\sigma$ puisque
$H \subset H'$.\footnote{La page écrit $\dot\sigma^{-1}H'\dot\sigma$ ; c'est le
même sous-groupe, puisque $\dot\sigma^{2} \in H \subset H'$.}

Inversement, toute donnée $(G, H \subset H', \dot\sigma)$ de ce type définit,
sur $S = G/H'$ et $A = G/H$, un graphe sur lequel $G$ agit transitivement
sur les sommets et sur les arêtes orientées --- ce que la page annonce par
« Si on a une telle situation » avant de s'interrompre. C'est la description
d'aujourd'hui des graphes \emph{arc-transitifs} comme graphes de classes
(\emph{coset graphs}), due pour l'essentiel à Sabidussi
(1964).\footnote{Rien sur la page ne cite Sabidussi, ni les travaux de Tutte
sur les graphes cubiques symétriques (1947).}

L'icosaèdre en est un exemple, et c'est une façon de relire le Théorème de la
page 7.\footnote{Ce rapprochement est le nôtre.} Pour $G = \Gamma' \simeq \mathfrak{A}_5$,
simplement transitif sur les arêtes orientées, on a $H = \{1\}$ ; $H'$ est le
groupe cyclique d'ordre $5$ des rotations autour de $u$ ; et $\dot\sigma$
est le demi-tour autour de l'axe qui passe par le milieu de l'arête
$\{u, x_0\}$, qui n'est pas dans $H'$. La condition simpliciale dit que les
stabilisateurs de deux sommets voisins ne se rencontrent qu'en l'identité.
On retrouve ainsi douze sommets et soixante arêtes orientées.

\end{document}
